% ===================================================================
% main.tex — 《经济学论文精读讲义》主文件
% 编译：XeLaTeX（中文 ctex 依赖；Overleaf 项目 Menu 中选 XeLaTeX，
%       TeX Live 版本按项目设置固定，见 CONTRIBUTING.md「编译与 Overleaf 对接」）
% 结构：\part = 经济学领域，\chapter = 一篇论文／一段领域史／一种计量方法；
%       每个 part 的第一章固定是 partinfo 章；每个 \part 前用 \partintro 写导语。
% ===================================================================
% openany（2026-10-21）：章/部分从任一页起，PDF 不出现整页留白；
% 付印装订时改回 openright（见 CONTRIBUTING「版式约定」）
\documentclass[openany]{ctexbook}

% ===================================================================
% preamble.tex — 经济学论文精读讲义公共导言
% 所有章节共享的宏包、TikZ 样式、自定义盒子与记号
% ===================================================================

% 版心: A4, 每行约 38 个汉字 (出版惯例 28--36, 讲义取上限以兼顾 A4 打印)。
% 付印装订时: 在下行加 bindingoffset=6mm (内侧留装订量), 并把 openany 改 openright。
\usepackage[a4paper, hmargin=3.0cm, vmargin=2.6cm]{geometry}
\usepackage{emptypage}   % openright 产生的空白左页不排页眉页码
\emergencystretch=1.5em  % 中西文混排偶发的断行困难兜底
\usepackage[round,authoryear]{natbib}   % 经济学惯例: 作者-年份夹注
\usepackage[font=small, labelfont=bf, skip=6pt]{caption}
\usepackage{imakeidx}   % 索引; 须在 hyperref 之前加载
\makeindex[intoc=true, title=索引]

% ------------------------- 版式约定 -------------------------
% 章标题沿用 ctexbook 默认（黑体大号居中）；节及以下顶格左对齐、黑体，
% 落实「黑体标题、宋体正文」的教材惯例（承旧仓 2026-08 定版，见 CONTRIBUTING）。
\ctexset{
  section       = {format = \Large\heiti\raggedright},
  subsection    = {format = \large\heiti\raggedright},
  subsubsection = {format = \heiti\raggedright},
  part/pagestyle = empty,   % 部分页不排页码
}
% 全书目录只列到节（x.y）；各章更深的导航由「导读地图」承担
\setcounter{tocdepth}{1}
% 页眉不将英文转全大写（book 类默认的 \MakeUppercase 在中文页眉里突兀）
\renewcommand{\chaptermark}[1]{\markboth{\CTEXthechapter\hspace{1em}#1}{}}
\renewcommand{\sectionmark}[1]{\markright{\thesection\hspace{1em}#1}}
% 部分页导语: 在 \part 之前调用 \partintro{...}, 导语会排在该部分页标题下方
\newcommand{\partintrotext}{}
\newcommand{\partintro}[1]{\gdef\partintrotext{#1}}
\ctexset{part/aftertitle = {\par\nobreak\vspace{5\baselineskip}
  {\centering\begin{minipage}{0.72\textwidth}
     \normalfont\zihao{-4}\setlength{\parindent}{2em}\partintrotext
   \end{minipage}\par}
  \global\let\partintrotext\relax\gdef\partintrotext{}}}
\usepackage{amsmath, amssymb, bm}
\usepackage{booktabs}
\usepackage{array}
\usepackage{multirow}
\usepackage{makecell}
\usepackage{xcolor}
\usepackage{tikz}
\usetikzlibrary{positioning, arrows.meta, fit, backgrounds, calc, shapes.misc, decorations.pathreplacing}
\usepackage{pgfplots}
\pgfplotsset{compat=1.18}
\usepackage[most]{tcolorbox}
\usepackage{listings}
% 代码样式 (贴 Stata/R/Python 代码、命令与配置用)。注意: listings 对中文支持差, 代码内注释请用英文。
\lstset{basicstyle=\small\ttfamily, backgroundcolor=\color{black!4},
  frame=single, rulecolor=\color{black!25}, breaklines=true,
  columns=flexible, keepspaces=true, showstringspaces=false,
  keywordstyle=\color{blue!60!black}, commentstyle=\color{green!40!black},
  stringstyle=\color{red!50!black}, xleftmargin=1em, xrightmargin=1em}
% 章末参考文献: 每章末尾只列本章引用的文献; globalcitecopy 使其同时汇入书末总表
\usepackage[globalcitecopy]{bibunits}
\usepackage{xurl}      % 允许长 URL 在任意字符处折行 (参考文献里的长链接会撑破版心)
\usepackage[colorlinks=true, linkcolor=blue!60!black, citecolor=blue!60!black, urlcolor=blue!50!black]{hyperref}
\hypersetup{pdftitle={经济学论文精读讲义}, pdfauthor={Sagewood AI Lab, IAH AI Group, Renmin University of China}, pdfsubject={梳理经济学各领域研究脉络：研究方法的发展与重点研究成果的发现}}
% 孤行寡行惩罚 (段首行落页尾/段尾行落页首)
\widowpenalty=8000
\clubpenalty=8000

% ------------------------- 颜色体系 -------------------------
% 配色取中性语义色, 含义随图赋值 (主线/分支线、方法/发现、处理/对照),
% 同一张图内同一含义固定同一色; 不预设某种图的专用色。
\definecolor{cMain}{RGB}{31,119,180}     % 主线/主对象 蓝
\definecolor{cBranch}{RGB}{255,127,14}   % 分支/被放弃的路线 橙
\definecolor{cFind}{RGB}{68,138,74}      % 发现/结论 绿
\definecolor{cAlt}{RGB}{148,103,189}     % 对照/替代路线 紫
\definecolor{cGray}{RGB}{238,238,238}

% ------------------------- TikZ 样式 -------------------------
% 常用图形: partinfo 的发展脉络图与时间线、识别策略示意、简单 DAG。
\tikzset{
  blk/.style={draw=black!60, rounded corners=2pt, minimum width=2.7cm,
              minimum height=0.8cm, font=\small, align=center, inner sep=2pt},
  mainblk/.style={blk, fill=cMain!15},     % 主线上的对象
  branchblk/.style={blk, fill=cBranch!18}, % 被放弃/被搁置的路线
  findblk/.style={blk, fill=cFind!14},     % 发现/结论
  altblk/.style={blk, fill=cAlt!12},       % 对照/替代路线
  opblk/.style={blk, fill=cGray, minimum width=1.9cm},
  ar/.style={-{Stealth[length=2mm]}, thick},
  dar/.style={-{Stealth[length=2mm]}, thick, dashed},  % 弱传承/被搁置的方向
  lbl/.style={font=\footnotesize\itshape, text=black!70},
  yr/.style={font=\small\bfseries, text=black!60},     % 时间线上的年份
}

% ------------------------- 自定义盒子 -------------------------
% 三类盒子按章编号 (识别要点 2.1、易错点 2.3 …), 三类计数器各自独立。
% 需被跨章引用的盒子, 在盒内首行放 \label{pit:xxx}, 引用处写「第 2 章的易错点~\ref{pit:xxx}」。
% 识别要点 (可整体搬用的模型设定、口径与识别设计)
\newtcolorbox[auto counter, number within=chapter]{ident}[1][]{colback=green!4,
  colframe=green!45!black, title={\textbf{识别要点~\thetcbcounter}~#1},
  fonttitle=\small, left=2mm, right=2mm, top=1mm, bottom=1mm, breakable}
% 易错点
\newtcolorbox[auto counter, number within=chapter]{pitfall}[1][]{colback=red!4,
  colframe=red!60!black, title={\textbf{易错点~\thetcbcounter}~#1},
  fonttitle=\small, left=2mm, right=2mm, top=1mm, bottom=1mm, breakable}
% 直觉/为什么
\newtcolorbox[auto counter, number within=chapter]{intuition}[1][]{colback=blue!4,
  colframe=blue!50!black, title={\textbf{直觉~\thetcbcounter}~#1},
  fonttitle=\small, left=2mm, right=2mm, top=1mm, bottom=1mm, breakable}
% 原文对照(备用: 体例是页码/表号随正文与盒标题就近标注, 本盒一般不用)
\newtcolorbox{origin}[1][]{colback=black!3, colframe=black!40,
  title={\textbf{原文出处}~#1}, fonttitle=\small, left=2mm, right=2mm,
  top=1mm, bottom=1mm, breakable}

% 章末参考文献: 每章正文末尾调用 \chapterrefs, 只列本章引用的文献
\defaultbibliographystyle{plainnat}   % 各章参考文献单元沿用与书末总表相同的样式
\newcommand{\chapterrefs}{%
  \putbib[bibs/refs]%
}
% 路径相对 main.tex (Overleaf 项目根 = knowledge/ 文件夹); 本地在 knowledge/ 内编译时相同。
% 章末书目（bibunits）用「本章参考文献」；书末总表另用「参考文献」，见 main.tex 的 \bibname 设置
\renewcommand{\bibsection}{\section*{本章参考文献}\markright{本章参考文献}}

% 跨章引用: \chref{ch:xxx} -> 「第1章」(无多余空格)
\newcommand{\chref}[1]{第\ref{#1}章}

% 章内建档/更新日期 (2026-10-22 加): 每章 \chapter 标题下方渲染一行小字。
% 用法: \chdate{2026-10-18}{2026-10-21}  (Created-At, Updated-At)。
% 注意: 不可用 \time{...} —— \time 是 TeX 只读寄存器(距午夜的分钟数), 不能带参数。
\newcommand{\chdate}[2]{{%
  \centering\noindent\footnotesize\color{black!55}
  Created-At: #1,\quad Updated-At: #2\par}}

% ------------------------- 实验室标识 -------------------------
% Sagewood AI Lab 主标记: 一片叶, 叶脉画成神经网络。改动请与 brand/ 下的 SVG 保持同步。
% 坐标沿用 SVG 的 64x64 画布(已翻转 y 轴), 自然尺寸约 22.5mm 高;
% 换尺寸用可选参数缩放, 如 \sagewoodmark[0.5]。
\definecolor{sagewood}{HTML}{4F7A5B}
\definecolor{sagewoodink}{HTML}{20342A}
\newcommand{\sagewoodmark}[1][1]{%
  \scalebox{#1}{%
  \begin{tikzpicture}[x=1pt, y=1pt,
      swleaf/.style={draw=sagewood, line width=2.6pt, line cap=round, line join=round}]
    \path[swleaf] (32,58) .. controls (46.5,45.5) and (50,26) .. (32,6)
                          .. controls (14,26)   and (17.5,45.5) .. (32,58) -- cycle;
    \path[swleaf] (32,52) -- (32,10);
    \path[swleaf] (32,19) -- (24.5,24.5);   \path[swleaf] (32,19) -- (39.5,24.5);
    \path[swleaf] (32,28) -- (24,33.5);     \path[swleaf] (32,28) -- (40,33.5);
    \path[swleaf] (32,36.5) -- (26,42);     \path[swleaf] (32,36.5) -- (38,42);
    \foreach \p in {(32,52),(24.5,24.5),(39.5,24.5),(24,33.5),
                    (40,33.5),(26,42),(38,42),(32,10)}
      { \fill[sagewood] \p circle (2.9pt); }
  \end{tikzpicture}}%
}

% 纯文字字标: 只有 SAGEWOOD / AI LAB 两行, 不含叶形标记。
% 用于同一页已出现过标记的场合(如扉页顶部已有标记时的底部署名),
% 避免同一符号在一页上重复出现。
\newcommand{\sagewoodwordmark}[1][1]{%
  \scalebox{#1}{%
  \begin{tabular}{@{}c@{}}
    {\color{sagewoodink}\fontsize{15.5}{15.5}\selectfont S\,A\,G\,E\,W\,O\,O\,D}\\[3pt]
    {\color{sagewood}\footnotesize A\,I\ \,L\,A\,B}
  \end{tabular}}%
}

% 横向组合标识: 标记 + SAGEWOOD + AI LAB。
% 对应 brand/sagewood-lockup.svg, 但字体随讲义正文, 不会有字体回退问题。
% 用于扉页署名; 也可用于章末、封底。缩放同上, 如 \sagewoodlockup[0.8]。
\newcommand{\sagewoodlockup}[1][1]{%
  \scalebox{#1}{%
  \begin{tikzpicture}[x=1pt, y=1pt]
    \node[anchor=west, inner sep=0] (swm) at (0,0) {\sagewoodmark[0.5]};
    \draw[sagewood!35, line width=0.5pt]
      ($(swm.east)+(8,-14)$) -- ($(swm.east)+(8,14)$);
    \node[anchor=west, inner sep=0] at ($(swm.east)+(17,6.5)$)
      {\color{sagewoodink}\fontsize{15.5}{15.5}\selectfont S\,A\,G\,E\,W\,O\,O\,D};
    % 中文名暂不启用; 需要时把下一行注释掉、启用再下一行即可
    \node[anchor=west, inner sep=0] at ($(swm.east)+(17,-8.5)$)
      {\color{sagewood}\footnotesize A\,I\ \,L\,A\,B};
    %\node[anchor=west, inner sep=0] at ($(swm.east)+(17,-8.5)$)
    %  {\color{sagewood}\footnotesize A\,I\ \,L\,A\,B\ \ %
    %   \color{black!45}· 智林人工智能实验室};
  \end{tikzpicture}}%
}

% ------------------------- 封面与 bjr 标识（2026-11-19 加） -------------------------
% 人大红：取通行的近似色值 #AE0B2A（RGB 174,11,42）。若拿到校方 VI 手册的标准色值，
% 只改下面这一行，标识与封面随之更新。
\definecolor{rucred}{HTML}{AE0B2A}
\definecolor{coverink}{HTML}{2B2B2B}

% bjr 标识：圆角方底上三个单线小写字母 b、j、r；j 的点同时是一条向下倾斜的需求曲线与
% 一条向上倾斜的供给曲线的交点（均衡点），取「经济学」之意。本标识为讲义自行设计，
% 只取人大红这一颜色，不含也不仿照任何学校徽志。
% 坐标为 64x64 画布（单位 pt）；全部用路径绘制，不依赖字体。
% 用法：\bjrmark 或 \bjrmark[0.6]（红底白字）；\bjrmarkrev[0.6]（白底红字，放在红色底上用）。
\newcommand{\bjrmarkdraw}[3]{% #1 缩放, #2 底色, #3 字色
  \scalebox{#1}{%
  \begin{tikzpicture}[x=1pt, y=1pt,
      bjrglyph/.style={draw=#3, line width=4.6pt, line cap=round, line join=round},
      bjrcurve/.style={draw=#3, line width=1.8pt, line cap=round, opacity=0.5}]
    \fill[#2, rounded corners=13pt] (0,0) rectangle (64,64);
    % 需求曲线（左上到右下）与供给曲线（左下到右上），交于 j 的点 (38,43)
    \path[bjrcurve] (23,57) .. controls (27,50) and (32,45.5) .. (38,43)
                            .. controls (43,41) and (47,40) .. (51,39.2);
    \path[bjrcurve] (26,39) .. controls (30,39.6) and (34,41) .. (38,43)
                            .. controls (44,46.5) and (50,51.5) .. (55,57.5);
    % b：竖笔加圆腹
    \path[bjrglyph] (12,17) -- (12,48);
    \path[bjrglyph] (20.5,25.5) circle (8.5);
    % j：竖笔带左钩；点即均衡点
    \path[bjrglyph] (38,34) -- (38,14) arc[start angle=0, end angle=-90, radius=6.5];
    \fill[#3] (38,43) circle (3.6pt);
    % r：竖笔加肩
    \path[bjrglyph] (48,34) -- (48,17);
    \path[bjrglyph] (48,26) .. controls (48,31.5) and (52,34.5) .. (56.5,33);
  \end{tikzpicture}}%
}
\newcommand{\bjrmark}[1][1]{\bjrmarkdraw{#1}{rucred}{white}}
\newcommand{\bjrmarkrev}[1][1]{\bjrmarkdraw{#1}{white}{rucred}}

% 封面：上部人大红色块（标识、书名），色块右侧以淡白线画 long-run 包络线与三条
% short-run 曲线（本讲义第一部分的核心图，也是微观经济学里长期与短期曲线的经典画法）；
% 下部白底放副题与署名。用 remember picture + overlay 定位，需编译两遍（latexmk 与
% Overleaf 均自动处理）。在 main.tex 的 \frontmatter 之后调用 \bjrcover。
\newcommand{\bjrcover}{%
  \begin{titlepage}
  \thispagestyle{empty}
  \begin{tikzpicture}[remember picture, overlay]
    \coordinate (blockSW) at ([yshift=-0.43\paperheight]current page.north west);
    \coordinate (blockSE) at ([yshift=-0.43\paperheight]current page.north east);
    \fill[rucred] (current page.north west) rectangle (blockSE);
    % 淡白线：long-run 包络线与三条在各自 climate 处相切的 short-run 曲线
    \begin{scope}
      \clip (current page.north west) rectangle (blockSE);
      \begin{scope}[shift={([xshift=0.45\paperwidth, yshift=58mm]blockSW)},
                    x=0.115\paperwidth, y=13mm, white, line cap=round]
        \draw[line width=2.4pt, opacity=0.22]
          plot[domain=0:4, samples=80] (\x, {4*\x-\x*\x});
        \foreach \m in {1,2,3}
          \draw[line width=1.3pt, opacity=0.22, dash pattern=on 5pt off 4pt]
            plot[domain={\m-1.05}:{\m+1.05}, samples=50] (\x, {4*\x-\x*\x-(\x-\m)*(\x-\m)});
        \foreach \m/\y in {1/3, 2/4, 3/3} \fill[opacity=0.38] (\m,\y) circle (2.6pt);
      \end{scope}
    \end{scope}
    % 标识（白底红字）
    \node[anchor=north west, inner sep=0] at ([xshift=28mm, yshift=-27mm]current page.north west)
      {\bjrmarkrev[0.95]};
    % 书名与英文副题
    \node[anchor=south west, inner sep=0, text=white] at ([xshift=27.2mm, yshift=24mm]blockSW)
      {\heiti\fontsize{37}{44}\selectfont 经济学论文精读讲义};
    \node[anchor=south west, inner sep=0, text=white, opacity=0.88] at ([xshift=28mm, yshift=14mm]blockSW)
      {\sffamily\bfseries\fontsize{9.5}{12}\selectfont
       E\,C\,O\,N\,O\,M\,I\,C\ \ \,L\,I\,T\,E\,R\,A\,T\,U\,R\,E\ \ \,·\ \ \,R\,E\,A\,D\,I\,N\,G\ \ \,N\,O\,T\,E\,S};
    % 下部白底：副题
    \node[anchor=north west, inner sep=0, text=coverink, align=left] at ([xshift=28mm, yshift=-17mm]blockSW)
      {\songti\zihao{4} 按领域梳理研究脉络：方法的发展与重点发现\\[7pt]
       \zihao{-4}\color{black!55} 讲清识别\quad·\quad 讲清取舍\quad·\quad 讲清测度};
    % 署名
    \draw[rucred, line width=1.6pt] ([xshift=28mm, yshift=47mm]current page.south west) -- ++(24mm,0);
    \node[anchor=north west, inner sep=0, text=coverink, align=left] at ([xshift=28mm, yshift=41mm]current page.south west)
      {\sffamily\bfseries\zihao{3} bjr\\[5pt]
       \songti\zihao{-4}\color{black!55} Sagewood AI Lab\quad·\quad \the\year~年~\the\month~月};
  \end{tikzpicture}
  \end{titlepage}%
}

% ------------------------- 无空白页（openany 下） -------------------------
% book 类的 \frontmatter 与 \mainmatter 无条件执行 \cleardoublepage，即使用了 openany，
% 目录止于奇数页时也会在正文前插入一张空白页。未用 openright 时改为 \clearpage；
% 付印改回 openright 后本段自动失效，恢复标准行为。
\makeatletter
\if@openright\else
  \renewcommand{\frontmatter}{\clearpage\@mainmatterfalse\pagenumbering{roman}}
  \renewcommand{\mainmatter}{\clearpage\@mainmattertrue\pagenumbering{arabic}}
\fi
\makeatother

% ------------------------- 数学记号 -------------------------
% 全书级记号在此登记; 各章可沿用所读论文的原始记号(首用处说明),
% 新造全书级记号前先查本块避免冲突, 新约定要登记进这里。
\newcommand{\R}{\mathbb{R}}
\newcommand{\E}{\mathbb{E}}                            % 期望
\newcommand{\Var}{\operatorname{Var}}
\newcommand{\Cov}{\operatorname{Cov}}
\newcommand{\Corr}{\operatorname{Corr}}
\newcommand{\ind}[1]{\mathbf{1}\!\left\{#1\right\}}    % 示性函数
\newcommand{\eps}{\varepsilon}                         % 误差项

% 书末总参考文献的标题（各章末文献表的标题「本章参考文献」在 preamble.tex 设置）
\renewcommand{\bibname}{参考文献}

\title{经济学论文精读讲义}

\author{bjr}

\begin{document}

% 按章分文献单元：各章末的 \chapterrefs（\putbib）只列本章引用
% （bibunits 要求在 \begin{document} 之后启用；见其手册 §3.2）
\bibliographyunit[\chapter]

\frontmatter
% 封面：人大红色块 + bjr 标识（宏在 preamble.tex「封面与 bjr 标识」）；
% 要回到 ctexbook 默认扉页，把下一行换回 \maketitle 即可
\bjrcover
\tableofcontents

\mainmatter

% ============================== 气候变化适应 ==============================
\partintro{本部分梳理 climate change adaptation 研究的脉络：climate 影响的估计方法如何
从截面 Ricardian 模型走到 weather 波动面板（panel approach），再到对「panel approach
能否反映 long-run adaptation」的系统检验。建议先读 partinfo 章把握脉络，再进各篇精读章
拆解 identification strategy。}
\part{气候变化适应}
\label{part:cca}

% partinfo 章：Climate Change Adaptation 领域的研究脉络（活页章）
% 建档/更新日期由章首 \chdate 渲染（grep: chdate）
% 维护纪律：每落地或大改一篇精读章，必须回到本章补定位、改叙事
%           （CONTRIBUTING「partinfo 维护纪律」）。
% 体例注记: 2026-10-21 起，正文中的 terms/expressions 一律用英文原词
%          （weather、climate、location、panel approach……），
%          中文仅保留句法成分；定义处随文给中文对照并登记英文 \index。
% 方法线核对（bjr 2026-10-18 提出，2026-11-04 扩展）：方法线按正典
%          papers/01_methods/ 下的原文逐站核对——deschenes2007economic、
%          mcintosh2006identifying（工作稿）、burke2015global、
%          hsiang2016climate（所据 NBER w22181 工作稿版）、merel2021climate，
%          关键主张照引原句（CONTRIBUTING「原句引用」2026-10-29 条）；
%          后续改方法线仍须回各正典 PDF 核对页码与版本。
% 方法线重构（2026-11-19）：按主笔 2026-11-18 的修改意见，方法线由原来的五站
%          单线改为两条线——cross-sectional line（cross-sectional approach →
%          long differences → multistage model）与 panel line（panel approach →
%          nonlinear panel approach → climate penalty approach），Hsiang (2016)
%          作为把两条线放进同一框架的一站单列。原五站的原句引用全部保留并逐条
%          对照本地 PDF 改回原文字面；新增 burke2016adaptation、da2024climate
%          两篇的原句各若干。multistage model 一站所据的 Butler 与 Huybers（2013）
%          与 Kolstad 与 Moore 的综述（所据 NBER w25537 工作稿版）于 2026-11-19 下载
%          入 papers/01_methods/，该站已据原文写出。

\chapter{Climate Change Adaptation 研究脉络梳理}
\label{ch:info-climate-change-adaptation}
\chdate{2026-10-18}{2026-11-19}

本 part 梳理 climate change adaptation 研究的脉络。领域围绕的问题只有一个：climate change 会带来多大损害，经济主体又能通过长期 adaptation 抵消其中多少。本领域随 climate change 议题的兴起而出现；随着 adaptation 成为研究重点，方法重心从「估计影响」移向「检验 adaptation」。发展史按「当时面对什么问题、用什么方法回答、得到了什么发现、遗留了什么问题」登记，被推翻或修正的结论照实收录。


本 part 的两个术语区分在这里先立起来（identification 上的含义随精读章展开）。第一，\index{weather}weather（天气）指年际与年内的波动，\index{climate}climate（气候）指 weather 的长期分布：用 weather 波动识别的是 short-run response（不含 long-run adaptation），long-run adaptation 要靠横截面 climatic differences 或 long differences 这类长期 variation（变异）识别，凡引用识别结论先分清 variation 来源。第二，\index{adaptation}adaptation（适应）指长期发生的调整（品种、灌溉、种植制度），\index{coping}coping（应对）指短期风险管理（借贷、储蓄、转移），行文不混用。

\section{时间线}

以年份为轴登记方法演进与代表性发现；写「待补」的格子是欠账，对应精读章落地即还（自查：\texttt{grep -n '待补' part/*/partinfo.tex}）。

\begin{itemize}
  \item \textbf{1994}：Mendelsohn、Nordhaus 与 Shaw 用美国县横截面上的农地价值对 climate 变量回归，开截面 \index{Ricardian model}Ricardian 路线 \citep{mendelsohn1994impact}。
  \item \textbf{2006}：McIntosh 与 Schlenker 的工作稿给出含 penalty 结构 $(x_{it}-\mu_i)^2$ 的 DGP 与「hybrid estimator」（全局平方项与组内平方项并设），并指出二次项使 group means 重新进入识别：``a quadratic explanatory variable causes group means to re-enter the identification''\citep[摘要页；模型见第~4 页、hybrid estimator 见第~5--6 页]{mcintosh2006identifying}。\chref{ch:merel-gammans-2021} 的模型、增广回归与 proposition 皆源于此。
  \item \textbf{2007}：Deschênes 与 Greenstone 把 weather 波动面板（panel approach）引入农业影响评估（county fixed effects 加年际 weather 波动），回应截面路线的 omitted variables 批评 \citep{deschenes2007economic}（\chref{ch:deschenes-greenstone-2007}）。
  \item \textbf{2013}：Butler 与 Huybers 逐县估计美国玉米单产对高温的敏感度，再用敏感度在各县之间的差异刻画 adaptation，是 multistage model（多阶段模型）的首例 \citep{butler2013adaptation}。
  \item \textbf{2015}：Burke、Hsiang 与 Miguel 用 166 国 1960--2010 年度面板估计温度对经济产出的 nonlinear effect，响应在年均 13~$^\circ$C 处最优 \citep{burke2015global}（\chref{ch:burke-hsiang-miguel-2015}）；二次设定是面板文献的常见选择（名单见 \citealp[第~1208 页脚注~3]{merel2021climate}，其中既有这篇论文，也有 Deschênes 与 Greenstone 2007）。
  \item \textbf{2016}：Burke 与 Emerick 用 \index{long differences}long differences（两个长时段均值之差），对比 long-run response 与 panel approach 给出的 short-run response，考察美国农业的 adaptation \citep{burke2016adaptation}。同年 Hsiang 的综述系统整理了 climate econometrics 的各种 research designs \citep{hsiang2016climate}（\chref{ch:hsiang-2016}）。
  \item \textbf{2021}：Mérel 与 Gammans 把「nonlinear panel approach 的系数在多大程度上反映 long-run adaptation」形式化为可计算的 asymptotic bias 分解问题，并提出 climate penalty 增广回归（\chref{ch:merel-gammans-2021}）。
  \item \textbf{2024}：Da、Xu、Yi 与 McCarl 把 climate penalty model 用于中国冬小麦县级面板，配合分季节的 degree-days 变量 \citep{da2024climate}（\chref{ch:da-et-al-2024}）。
  \item 待补：Schlenker 与 Roberts（2009）的非线性温度响应一站；结构外推一站；各点落位后本时间线与「发现线」互相对照修订。
\end{itemize}

\section{方法线}

方法沿两条线演进，划分的依据是用来识别 climate 效应的 variation 来自哪里。\textbf{cross-sectional line} 用 location 之间的 climatic differences：从最早的 cross-sectional approach，到 long differences 与 multistage model。\textbf{panel line} 用 location 内部逐年的 weather fluctuation：从最早的 panel approach，经 nonlinear panel approach，到 climate penalty approach。两条线面对的是同一组取舍：横截面 variation 含有 long-run adaptation，但容易混入 omitted variables；组内 variation 干净，但测到的是否含 adaptation 存疑。每一站都按「上一站解决不了什么问题，这一站用什么 variation 回答」登记；各站的关键主张尽量照引原文原句（CONTRIBUTING「原句引用」），页码为所据版本页码（Hsiang 2016 所据为 NBER 工作稿版，见本章参考文献注）。

\subsection{Cross-sectional line：用 location 之间的 climatic differences}

\textbf{第一站：cross-sectional approach（Ricardian，1994，见时间线）。}用横截面上的农地价值对 climate 变量回归。\index{cross-sectional approach}Cross-sectional approach（截面法）的长处，\citet{hsiang2016climate} 的概括是：长期居住于一地的农民深知本地 climate 并据此调整投资与管理，``farm prices can be assumed to reflect all direct effects as well as all belief effects''\citep[第~5 页；\index{direct effect}direct effects（直接效应）指 climate 经由 weather 实现的直接影响，\index{belief effect}belief effects（信念效应）指主体因相信本地 climate 如此而做的调整，即 adaptation 行为，区分随 Hsiang 2016]{hsiang2016climate}。其弱点是 omitted variables：``A weakness of the cross-sectional approach is its vulnerability to omitted variables bias''\citep[第~5 页]{hsiang2016climate}，且 ``an econometrician can never be certain their model is unbiased''\citep[第~5 页]{hsiang2016climate}。\citet{deschenes2007economic} 用数据证实了这一弱点：同一 hedonic 设定下，估计值对控制变量、样本与加权的微小选择极其敏感，``it produces estimates that are extremely sensitive to seemingly minor choices about control variables, sample, and weighting''\citep[第~354 页]{deschenes2007economic}。这一站留下的问题是：能不能既用到 climate 的长期差异，又摆脱 location 之间不可观测的水平差异。

\textbf{第二站：long differences（2016，见时间线）。}\citet{burke2016adaptation} 的做法是对每个县取两个长时段的均值之差，再跨县比较这些长期变化：``We adopt a `long differences' approach and model county-level changes in agricultural outcomes over time as a function of these changes in temperature and precipitation, accounting for time-invariant unobservables at the county level and time-trending unobservables at the state level''\citep[第~107 页]{burke2016adaptation}。它保留了横截面比较（各县变暖的幅度不同），但比较的是变化而不是水平，县级不随时间变化的不可观测因素因此被差分掉。\citet{hsiang2016climate} 给它的方法定位是折中：``An approach that aims to compromise between the strengths and weaknesses of cross-sectional analysis and time series identification is the long differences strategy, where changes for both the outcome and the climate within locations are correlated across locations''\citep[第~7 页]{hsiang2016climate}。long differences 给出的 long-run 估计可以与 panel approach 的估计对比，两者之差被读作 adaptation；Burke 与 Emerick 的结论见发现线。代价有两条：识别所用的是各 location 长期趋势之间的差异，样本里必须有足够多的 location 和足够大的趋势差异；作对比基准的面板估计本身是否纯粹是 short-run response，要到 panel line 第三站才有答案。

\textbf{第三站：\index{multistage model}multistage model（多阶段模型；Butler 与 Huybers 2013）。}这一站把「响应本身随 climate 而变」直接估出来，分两步做。\citet{kolstad2019estimating} 的综述这样概括：``A first step models the local-linear effect of weather variation, allowing this linear response to weather fluctuations to vary across space, all linearly. In a second step, the coefficient on weather (from the first step) is used as the left hand side variable with a nonlinear function of climate and other control variables on the right hand side''\citep[第~11 页，NBER w25537 版页码]{kolstad2019estimating}。第一步用各 location 内部的 weather 波动，第二步用 location 之间的 climatic differences。Kolstad 与 Moore 因此把它与 long differences 一并归入兼用两种 variation 的 hybrid approaches（第~6 页）；本讲义随 \citet{da2024climate} 把它放在 cross-sectional line 上，因为 long-run adaptation 的信息全部来自第二步的横截面比较。首例是 \citet{butler2013adaptation}：美国东部 1{,}600 多个县、1981--2008 年的玉米单产，温度的影响用 growing degree days 与 killing degree days（日最高温超过 29\,$^\circ$C 部分的累积）刻画；各地单产对后者的敏感度不同，较热的地方敏感度较低。把这种空间上的 adaptation 当作未来 adaptation 的替代，``we find that losses to average US maize yields from a 2~$^\circ$C warming would be reduced from 14\% to only 6\% and that loss in net production is wholly averted''\citep[第~68 页]{butler2013adaptation}。这一站的代价有两条。其一，第二步是截面回归，第一站的 omitted variables 问题在这里重现：``any variable correlated with climate that also affects the marginal response to weather will bias the estimate of the effect of climate if not explicitly controlled for''\citep[第~12 页]{kolstad2019estimating}。其二，它与 long differences 一样需要足够多的 location，\citet{da2024climate} 的说法是：``These approaches exploit cross-sectional variation and thus require a large number of counties to obtain reliable estimates''\citep[第~825 页]{da2024climate}。Butler 与 Huybers 一文本讲义只读了与设计有关的部分，未设精读章。

\subsection{Panel line：用 location 内部的 weather fluctuation}

\textbf{第一站：\index{panel approach}panel approach（面板法）的引入（2007，\chref{ch:deschenes-greenstone-2007}）。}\citet{deschenes2007economic} 的设计陈述是 ``This paper's idea is simple—we exploit the presumably random year-to-year variation in temperature and precipitation to estimate whether agricultural profits are higher or lower in years that are warmer and wetter''\citep[第~355 页]{deschenes2007economic}，估计 ``These data are used to estimate the effect of weather on agricultural profits and yields, conditional on county and state by year fixed effects''\citep[第~355 页；此思路本身可追溯到 Huntington 1922，考据见 \citealp{hsiang2016climate} 第~6--7 页]{deschenes2007economic}。它回应的正是 cross-sectional line 第一站的 omitted variables 批评：``The core benefit of this approach is that it accounts for unobservable differences between units, eliminating a potential source of omitted variables bias''\citep[第~7 页]{hsiang2016climate}。variation 来源换成逐年 \index{weather fluctuation}weather fluctuation（天气波动）后，两重局限随即显形，且都是 D-G 自己写明的。其一，climate 信息进不了模型：``Second, it is impossible to estimate the effect of the long-run climate averages in a model with county fixed effects, because there is no temporal variation in $\overline{W}_{ic}$''\citep[第~367 页；$\overline{W}_{ic}$ 为 climate normals，即各县的长期 climate 变量——\citet{merel2021climate} 第~1207 页所引即此句]{deschenes2007economic}。其二，面板估计到的是不含全部 substitution 的响应：``The import for the subsequent analysis is that our estimates of the impact of climate change may be downward biased, relative to the preferred long-run effect that allows for all economic substitutions''\citep[第~359 页]{deschenes2007economic}——这句自认是后来 adaptation 批评的源头。

\textbf{第二站：\index{nonlinear panel approach}nonlinear panel approach（非线性面板法，weather 变量以二次项等非线性形式进入面板）。}这一站的模型源头是 \citet{mcintosh2006identifying}（2006，工作稿）：其式 (1) 即含 penalty 结构 $(x_{it}-\mu_i)^2$ 的 DGP，其 ``hybrid estimator''（全局平方项与组内平方项并设）是 \chref{ch:merel-gammans-2021} 增广回归的前身，摘要一句 ``a quadratic explanatory variable causes group means to re-enter the identification''\citep[摘要页；式 (1) 在第~4 页，hybrid estimator 在第~5--6 页]{mcintosh2006identifying} 即 proposition 的最早表述；``climate penalty'' 这个名字则由 Mérel 与 Gammans 命名\citep[第~1213 页]{merel2021climate}。宏观尺度上的代表应用是 \chref{ch:burke-hsiang-miguel-2015}：\citet{burke2015global} 用 166 国 1960--2010 年度面板，``comparing a country to itself in years when it is exposed to warmer- versus cooler-than-average temperatures due to naturally occurring stochastic atmospheric changes''\citep[第~236 页]{burke2015global}，配二次温度项，得到 ``country-level economic production is smooth, non-linear, and concave in temperature (Fig.~2a), with a maximum at 13~$^\circ$C''\citep[第~236 页]{burke2015global}。这一步改变了 identification 的 variation 构成——非线性项与 \index{fixed effects}fixed effects 并用，使横截面 \index{climatic differences}climatic differences（气候差异）重新进入模型识别——由此形成 \citet{merel2021climate} 整理的 proposition（命题）：``Including nonlinearities in explanatory weather variables in a fixed-effects panel regression allows cross-sectional weather variation to enter model identification, suggesting that the estimates obtained from a nonlinear panel model at least partially reflect long-run adaptation to climate''\citep[第~1208 页]{merel2021climate}。该 proposition 在 2021 年前已被多项研究谨慎承认（名单见同页正文，其中有 \citealp{mcintosh2006identifying} 与 \citealp{burke2015global}），但它只回答了「估计里有没有 long-run 成分」：``the extent to which nonlinear panel estimates should be thought of as inclusive of climatic adaptation remains unclear''\citep[第~1208 页]{merel2021climate}。两点须防误读：其一，\citet{burke2015global} 是这一站的旗舰应用而非首创；其二，「线性对非线性」是教学简化——\citet{deschenes2007economic} 的利润方程已对 degree-days 与降水取二次（第~367 页），故它同列于 \citet{merel2021climate} 的二次设定面板名单（第~1208 页脚注~3）；真正决定系数如何解释的，是 penalty 结构是否在回归里。

\textbf{第三站：climate penalty approach（2021，\chref{ch:merel-gammans-2021}）。}这一站回答 proposition 悬置的「多少」：在 McIntosh 与 Schlenker（2006）提出的模型框架上（\citealp{mcintosh2006identifying}），把「nonlinear panel estimates 里 long-run 成分占多少」形式化为可计算的 \index{asymptotic bias}asymptotic bias（渐近偏差）分解。其一，``panel estimates from the quadratic weather-outcome relationship can be thought of as a weighted average of short- and long-run responses, with the weight on the long-run coefficients increasing with the share of overall weather variation attributable to cross-sectional (i.e., climatic) differences''\citep[第~1209 页]{merel2021climate}；weather 分布无 skewness 时为 convex combination（凸组合，第~1217 页），线性情形下偏差方向由条件于 location 的 weather 分布的 \index{skewness}skewness（偏度）决定，可正可负（第~1215 页）。其二，同一框架反向给出一致估计的路线：在面板中显式加入 \index{climate penalty}climate penalty term，一次回归同时得到 conditional on climate 的 short-run response、全局 long-run response 与一个 \index{climatic adaptation}climatic adaptation 的正式检验（第~1210、1220--1221 页）；该检验与 long differences 的检验同型，``our test of climatic adaptation is implicit in that it relies on the comparison of short-run versus long-run outcome responses''\citep[第~1210 页]{merel2021climate}。这一设定此后被直接搬用到新的作物与国家：本 part 的 Application 1（\chref{ch:da-et-al-2024}）把它用于中国冬小麦，配合分季节的 degree-days 变量，并报告了 penalty term 全部加入时出现的反常符号；它成为这条线上通行的做法。模型 terms 与解释 terms 在这条线上的逐条对应与显式分解，见精读章 \chref{ch:merel-gammans-2021}。

\subsection{把两条线放进同一框架：Hsiang（2016）}

\chref{ch:hsiang-2016} 把上述争论正式化。\citet{hsiang2016climate} 的出发点是 ``a debate as to whether regressions on `weather' variables provide meaningful insight into the effects of climate''\citep[第~2 页]{hsiang2016climate}，他给出 ``conditions under which the use of weather variables in regressions is justified and, perhaps surprisingly, dominates traditionally preferred methods''\citep[第~2 页]{hsiang2016climate}，并把领域整理成三种 research designs：``cross-sectional approaches, use of time-series variation, and a hybrid known as `long differences'''\citep[第~5 页]{hsiang2016climate}——前者与后者属 cross-sectional line，中间一种属 panel line。他提出两件工具。其一，\index{frequency-identification tradeoff}frequency-identification tradeoff（频率—识别权衡）：低频 variation 可能捕捉 adaptation，但识别可信度随之下降（第~8--9 页）；若 belief effects 主要通过改变 direct effects 起作用，``total effects of climate still may be nearly identified with high-frequency time-series''\citep[第~9 页]{hsiang2016climate}。其二，对 \index{marginal treatment comparability}marginal treatment comparability（边际处理可比性，第~6 页）的部分检验（第~9--10 页）：若三种设计的估计相等（$\hat\beta_{CS}=\hat\beta_{LD}=\hat\beta_{TS}$），weather 的边际变化即可当作 climate 边际变化的替代；他在玉米单产数据上实施的结果，除横截面外各估计统计上无异（第~10 页）。他还证明了一个正式结果：若结果变量已被优化，在边际处 weather 效应等于 climate 效应（\S 2.4）。与这一结果对照，Mérel 与 Gammans 自陈：``Whereas previous research has already argued that short-run and long-run responses should be identical at the margin whenever the outcome variable is optimized (e.g., Hsiang 2016), our result is both different and more specific''\citep[第~1217 页]{merel2021climate}——面板估计到的是加权平均，恰在边际处等同只是特例。

两条线各站的适用边界比较，以及与之平行的结构外推一站，待补。

\section{发现线}

待补：本领域已确认的重点结论与量级区间，随各精读章「结果与解释」节落地回填；数字必须带表号（页码）与样本口径。不同研究的 climate 变量定义、样本与被解释变量不同，下列数字不得跨条比较大小。目前已可登记六条。

其一与其二出自 \chref{ch:merel-gammans-2021}：摘要报告对美国与法国的作物单产找到了 long-run climatic adaptation 的一些证据；方法线悬置的「面板系数里 long-run 成分占多少」第一次有了量级回答——growing season 平均温度的效应应主要按 long-run response 解读，三个面板的长期权重分别为 98\%、88\%、86\%，growing season 降水的效应则含相当大的 short-run 成分，长期权重分别为 41\%、7\%、53\%（Table~2 与正文，第~1225 页；样本口径：美国县面板为 100 度经线以东、不含 Florida 的县，$I=2{,}391$；Corn Belt 三州 Iowa、Illinois、Indiana，$I=293$；法国 department 级面板，$I=88$；三面板均 $T=66$；美国的产出为玉米与大豆单产，法国为小麦与大麦单产，第~1224 页）。

其三出自 \chref{ch:deschenes-greenstone-2007}：``Using long-run climate change predictions, the preferred estimates indicate that climate change will lead to a \$1.3 billion (2002\$) or 4.0-percent increase in annual agricultural sector profits''\citep[第~356 页，AER 97(1) 已刊版口径]{deschenes2007economic}，且响应面在预测变化区间内平坦——``all of the response surfaces are flat over the ranges of the predicted changes in degree-days and inches''\citep[第~355 页]{deschenes2007economic}。

其四出自 \chref{ch:burke-hsiang-miguel-2015}：国家层面的经济产出对年平均温度呈光滑的凹响应，最优在 13~$^\circ$C（166 国、1960--2010 年、人均 GDP 增长率，第~236 页）；1960--1989 与 1990--2010 两段子样本的响应几乎相同，作者据此判断样本期内未见明显 adaptation（第~237 页）；基准 projection 下 2100 年全球产出比无 climate change 的情形低 23\%，影响为正的概率 0.29（SSP5，第~238 页）。

其五出自 Burke 与 Emerick（2016，本讲义未设精读章）：``Longer run adaptations appear to have mitigated less than half—and more likely none—of the large negative short-run impacts of extreme heat on productivity''\citep[第~106 页，摘要]{burke2016adaptation}（美国县级农业，long differences 对面板估计）。

其六出自 \chref{ch:da-et-al-2024}（中国黄淮海 352 县、1981--2015 年、冬小麦对数单产）：是否纳入 freezing 变量决定 projection 的符号，不计时各情景减产 12.5\% 到 48.2\%，计入时小幅增产 1.9\% 到 2.5\%（第~817、837 页）；long-run 口径下 Fall 高于 24\,$^\circ$C 的 degree-days 每多 10 个减产 7.0\%，conventional 面板的对应数字是 11.4\%（Table~2，第~829、832 页）。

待补：被推翻或修正的结论（如早期 Ricardian 外推的争议）。已可登记两例。第一例，\chref{ch:deschenes-greenstone-2007} 重做 hedonic 估计，``estimates of the effect of the benchmark doubling of greenhouse gases on the value of agricultural land range from $-\$200$ billion (2002\$) to \$320 billion (or $-18$ percent to 29 percent)''\citep[第~357 页]{deschenes2007economic}，比此前文献记录的区间更宽，结论是 ``Despite its theoretical appeal, we conclude that the hedonic method may be unreliable in this setting''\citep[第~357 页]{deschenes2007economic}。第二例是对一种做法的修正：把 long differences 估计与 conventional 面板估计之差当作 adaptation 的量度，前提是面板估计代表 short-run response；\chref{ch:merel-gammans-2021} 说明面板估计本身是 short-run 与 long-run response 的加权平均，\chref{ch:da-et-al-2024} 据此指出这种比较可能低估 adaptation（该文第~836 页）。

\section{章节速查表}

本表是「章与 partinfo 互指」的锚：每篇精读章一格。讲方法的各章在前，应用篇在后；应用篇按收入本 part 的先后编号（Application 1、2……），编号出现在目录、页眉、本表与各章的坐标句里，章题仍用论文标题。

\begin{table}[htbp]
\centering
\small
\begin{tabular}{@{}p{1.4cm}p{2.8cm}p{3.0cm}p{2.7cm}p{3.2cm}@{}}
\toprule
章 & 论文 & 回答的问题 & 方法 & 给领域留下了什么 \\
\midrule
\chref{ch:deschenes-greenstone-2007} & The Economic Impacts of Climate Change: Evidence from Agricultural Output and Random Fluctuations in Weather（AER 2007） & 年际 weather 波动识别的 climate change 对美国农业利润的影响有多大？hedonic 截面估计是否可靠？ & county 与 state$\times$year fixed effects 下的年际 weather 波动回归（panel line 第一站） & hedonic 估计对口径选择极敏感、不可靠的重估证据；自认面板估计不含全部 substitution、相对 long-run effect 向下偏（adaptation 批评的源头） \\
\chref{ch:burke-hsiang-miguel-2015} & Global Non-Linear Effect of Temperature on Economic Production（Nature 2015） & 温度与经济产出在全球国家面板上是否非线性？样本期内的变暖是否带来了 adaptation？ & 166 国 1960--2010 增长率面板，二次温度项（country 与 year fixed effects、国家二次趋势；panel line 第二站的旗舰应用） & 全球温度—产出响应在 13~$^\circ$C 处最优的非线性证据；「样本期内未见 adaptation」的判断引发后续 adaptation 检验 \\
\chref{ch:hsiang-2016} & Climate Econometrics（ARRE 2016） & weather 回归何时能识别 climate 效应？三种 research designs 如何取舍？ & 方法综述：正式化 weather／climate，把 cross-section、time-series、long differences 三种设计放进同一框架 & marginal treatment comparability、frequency-identification tradeoff；$\hat\beta_{CS}=\hat\beta_{LD}=\hat\beta_{TS}$ 部分检验（玉米单产应用）；边际处 weather 效应等于 climate 效应的正式结果 \\
\chref{ch:merel-gammans-2021} & Climate Econometrics: Can the Panel Approach Account for Long-Run Adaptation?（AJAE 2021） & nonlinear fixed effects 面板（忽略 penalty term 的二次设定）估计到的系数，在多大程度上能当作 long-run climate response 来读？ & 在 McIntosh 与 Schlenker（2006）框架上做 asymptotic bias 分解（固定 location 数、大 $T$）；climate penalty 增广回归同时估计 long-run 与 short-run response（panel line 第三站） & 线性情形偏差由条件于 location 的 weather skewness 驱动、可正可负；无 skewness 时面板系数是 short-run 与 long-run response 的 convex combination；空间大面板的温度系数已主要反映 long-run response（降水不然） \\
\chref{ch:da-et-al-2024}（Application 1） & Climate Change and Winter Wheat Yields: The Roles of Freezing Days and Long-Run Adaptations（ERAE 2024） & 中国冬小麦的 climate 影响评估是否必须分季节并纳入 freezing？climate penalty model 能否用于区分 short-run 与 long-run 影响？ & 黄淮海 352 县 1981--2015 面板；按季节由数据选阈值构造 degree-days；对两个变量加 penalty term（panel line 第三站的应用） & freezing 变量决定 projection 的符号；penalty term 全部加入时出现反常符号的实例；long differences 与面板之差可能低估 adaptation 的论点 \\
\bottomrule
\end{tabular}
\end{table}

% 章末参考文献（只列本章引用的条目），勿删
\chapterrefs

% ---- 方法篇：讲 research design 与 identification 的各章 ----
% 精读章: Deschênes & Greenstone, The Economic Impacts of Climate Change:
%         Evidence from Agricultural Output and Random Fluctuations in Weather
%         (AER 2007，97(1) 三月号正刊)
% 原文出处: doi:10.1257/aer.97.1.354（所据版本：期刊已刊版 AER 97(1): 354--385；
%          本仓 papers/01_methods/2007-AER-Deschenes-Greenstone-*.pdf 为团队正典，
%          本章页码引注均指该版）
% 体例注记: 2026-10-21 起，正文中的 terms/expressions 一律用英文原词，
%          中文仅保留句法成分；定义处随文给中文对照并登记英文 \index。
% 建档/更新日期由章首 \chdate 渲染（grep: chdate）
% 主笔: bjr；初稿: 2026-11-04（Claude 协助起草，bjr 审订）
% 核对: 2026-11-19 对照本地 PDF 全章重核——原句引用逐条改回原文字面；
%       改正 Table 2C、Table 3、SHF 复刻、Table 5/6/8、样本县数、置信区间、
%       分州显著性等处的数字与归属；式 (3)(4) 记号改从原文。
% partinfo 定位: 把 weather 波动面板（panel approach）引入农业影响评估，
%               并自认面板估计不含全部 substitution、相对 long-run effect
%               向下偏（adaptation 批评的源头），见
%               \chref{ch:info-climate-change-adaptation} 的方法线与速查表。

\chapter{The Economic Impacts of Climate Change（Deschênes \& Greenstone，AER 2007）}
\label{ch:deschenes-greenstone-2007}
% 章题过长，页眉用短题（preamble 的 \chaptermark 经 \markboth 排页眉）
\chaptermark{Economic Impacts of Climate Change（Deschênes \& Greenstone，2007）}
\chdate{2026-11-04}{2026-11-19}

\section{这篇论文在说什么}

一句话概括：论文用 counties 内部逐年 \index{weather}weather 波动（conditional on county 与 state$\times$year fixed effects）估计 climate change 对美国农业部门利润的影响，以此回应截面 hedonic（即 \index{Ricardian model}Ricardian）路线的 omitted variables 批评；其经验结论是影响小且不显著——``The preferred estimates indicate that climate change will increase annual profits by \$1.3 billion in 2002 dollars (2002\$) or 4 percent''\citep[第~354 页]{deschenes2007economic}，并判定截面路线``extremely sensitive to seemingly minor choices about control variables, sample, and weighting''\citep[第~354 页]{deschenes2007economic}，故不可靠。

它在 partinfo 里的坐标：panel approach 的引入站（\chref{ch:info-climate-change-adaptation} 方法线 panel 线第一站）。论文自己写明的局限——面板估计不含全部 adaptation、相对 long-run effect 向下偏（第~359 页）——就是后来 adaptation 检验路线的起点，由 \chref{ch:burke-hsiang-miguel-2015} 与 \chref{ch:merel-gammans-2021} 接手。

骨架表（所据版本：AER 97(1) 已刊版；页码为期刊页码）：

\begin{table}[htbp]
\centering\small
\begin{tabular}{@{}p{3.4cm}p{1.6cm}p{7.8cm}@{}}
\toprule
原文章节 & 页码 & 讲什么 \\
\midrule
Abstract & 354 & +\$1.3 billion（4\%）的基准结论；hedonic 不可靠 \\
引言（未编号） & 354--357 & production function、hedonic 两条旧路线的缺陷；论文设计与主要发现 \\
\S I Conceptual Framework & 357--360 & I.A 面板新策略与 hedonic 的差别；I.B 用年际波动推断 climate 效应的经济学（式 (1)(2)、Figure~2），自认 downward bias \\
\S II Data Sources and Summary Statistics & 360--365 & Census of Agriculture 县级面板；PRISM 与气象站数据、growing season degree-days 构造；climate 预测情景；weather variation 统计（Table~2） \\
\S III Econometric Strategy & 365--367 & III.A hedonic 设定（式 (3)）；III.B 面板设定（式 (4)）与 identification assumption \\
\S IV Results & 367--380 & IV.A hedonic 重估（Table~3--4）；IV.B 利润面板（Table~5--7）；IV.C 单产（Table~8） \\
\S V Interpretation & 380--381 & 福利换算（现值 \$26 billion）；四条 caveats，另加两条会使损害被高估的 caveats \\
\S VI Conclusions & 381 & 结论与后续研究方向 \\
\bottomrule
\end{tabular}
\end{table}

\section{预备知识}

读者需备两件旧工具。其一，论文之前的两条路线各有一难：\index{production function approach}production function approach（生产函数法）估计 weather 对具体作物单产的影响，剔除了土壤等农户控制不了的因素，但``farmer adaptations are completely constrained in the production function approach''\citep[第~355 页]{deschenes2007economic}，估计向下偏；\index{hedonic approach}hedonic approach（特征价格法，即 Ricardian 路线）用农地价值对 climate 变量回归，理论上地价``reflect the present discounted value of land rents into the infinite future''\citep[第~355 页]{deschenes2007economic}，``In principle, this approach accounts for the full range of farmer adaptations''\citep[第~355 页]{deschenes2007economic}，但 unmeasured characteristics（土壤质量、转作他用的期权价值等，作者引 Hoch 1958、1962 与 Mundlak 1961）与 climate 相关，``the sign and magnitude of the resulting omitted variables bias is unknown''\citep[第~355 页]{deschenes2007economic}。其二，\index{degree-days}degree-days（度日，又称生长度日；术语来自农学与气象学）：作物累积热量暴露的度量，日平均温度在基点与顶点之间时按超出基点的度数累积（构造口径见下文数据节）。

\section{数据与测度}

\textbf{产出与利润。}农业数据来自 Census of Agriculture（1978、1982、1987、1992、1997、2002 六次普查），观测单位是县；普查的填报口径是当年农产品销售额不低于 \$1{,}000 的农场与牧场（第~360 页）。生产支出在 1978 与 1982 年没有采集，所以利润分析以 1987--2002 四次普查为基础（第~360 页）。样本有三个口径，引用时要分清：Table~1 的描述统计用 1978--2002 的平衡面板，2{,}268 县（第~362 页及脚注~11；Alaska 与 Hawaii 剔除，土壤、农地面积、灌溉面积、人均收入、人口密度、纬度任一变量缺失的县剔除）；利润回归再要求利润变量在各普查年都不缺失，得 2{,}262 县的平衡样本（第~372 页脚注~23）；hedonic 重估（Table~4）的样本又少 144 个县（第~371 页）。被解释变量是县级每英亩农地的利润：农产品销售的市场价值减去生产支出，\textbf{不含联邦农业项目支付}（第~360 页）。作者明言这是不理想的口径：普查的支出项不含未抵押自有资本的利息，只部分计入了农场的资本成本（第~360 页），利润因而偏高；全县平均地租约 \$35（2002\$），利润约 \$42（2002\$），因此把利润口径的影响估计乘以二者之比 0.83 才换算成租金口径的福利度量（第~360 页）。

\textbf{climate 与 weather。}climate 数据用 PRISM 模型（4$\times$4 公里格点的降水与温度插值，作者称其被``almost all professional weather services... regarded as one of the most reliable interpolation procedures''\citep[第~361 页]{deschenes2007economic}）。growing season（生长季）取 4 月 1 日至 9 月 30 日，覆盖冬小麦以外的多数作物（第~361 页）。degree-days 的构造遵循农学传统（Ritchie 与 NeSmith 1991）：基点 46.4\,$^\circ$F（8\,$^\circ$C）、顶点 89.6\,$^\circ$F（32\,$^\circ$C），日平均温度在区间内贡献（均温$-$基点）个 degree-days，超过顶点的日子贡献 43.2 个 degree-days（第~361--362 页）。Ritchie 与 NeSmith 还讨论过 93.2\,$^\circ$F（34\,$^\circ$C）处可能存在一个有害温度阈值，超过它升温即造成损害；作者在 robustness check 里考察了这一点（第~361 页；即 Table~6 的 harmful degree-days 一行，第~376 页）。月度 PRISM 数据算不出逐日 degree-days，作者改用约 8{,}000 个地面气象站的逐日数据（NCDC Cooperative Summary of the Day）——这是对既有文献（用月度数据加分布假设，如 Thom 1966 与 SHF 2006）的测度改进（第~362 页）。

\textbf{climate 情景。}两组预测：基准情景（全季全区均匀升温 5\,$^\circ$F、降水 +8\%）；Hadley 2 模型的逐日州级预测（2020--2049 与 2070--2099 两窗口，第~362 页）。

\textbf{weather 波动与 climate change 的量级对比。}这是测度一节最该记住的表。Table~2C（第~364 页）报告去年份 fixed effects、去 state$\times$year fixed effects 之后，县年观测的偏离超过各幅度的占比（1987、1992、1997、2002 四年平均）。去 state$\times$year 效应后，24.5\%、9.3\%、2.2\% 的县年 degree-days 偏离超过 400、800、1{,}200；62.3\%、35.3\%、18.1\% 的县年 growing season 降水偏离超过 1.0、2.0、3.0 英寸（第~365 页）。把它与要外推的幅度对比：Hadley 2 长期情景的平均增幅约 1{,}200 degree-days（Table~2B；基准情景的增幅约 720 degree-days，Table~2A），这样大的 degree-days 偏离在数据里只来自 2.2\% 的观测，作者的措辞相应有所保留（幅度约 1{,}000 degree-days 的偏离占 5\%）；降水则宽裕得多，超过 18\% 的观测偏离大于 3.0 英寸（第~365 页）。结论是 identification 所用的 variation 与要外推的 variation 在量级上有重叠，降水的重叠充分，degree-days 的重叠落在分布的尾部。这是 panel approach 可用于外推的前提，也是它的边界。

\section{模型设定与识别策略}

本节是论文的理论贡献：先把截面 hedonic 写出来（式 (3)），再给出面板替代（式 (4)），由 identification 反推 assumption，最后用 \S I.B 的经济学说明这套估计量测到的是什么、偏向哪边。记号随原文：$\overline{W}_{ic}$ 是县 $c$ 的第 $i$ 个 climate 变量（长期均值，原文称 climate normals），$W_{ict}$ 是它在年 $t$ 的 weather 实现。

\subsection{截面 hedonic 及其问题}

hedonic 设定为（式 (3)，第~365 页）
\begin{equation}
y_{ct} = \mathbf{X}_{ct}'\boldsymbol{\beta} + \sum_i \theta_i f_i(\overline{W}_{ic}) + \epsilon_{ct},
\qquad \epsilon_{ct} = \alpha_c + u_{ct},
\label{eq:dg-hedonic}
\end{equation}
其中 $y_{ct}$ 是县 $c$ 在年 $t$ 的每英亩农地价值；climate 变量取 1 月、4 月、7 月、10 月的温度与降水共八个，这一选择沿 Mendelsohn、Nordhaus 与 Shaw（1994）；$f_i$ 取线性加二次项是文献惯例，作者写明各 climate 变量的恰当函数形式并不清楚（第~366 页）；灌溉县与非灌溉县参数分开，则是采纳 SHF（2005）的意见（第~366 页）。一致估计要求 $\E[f_i(\overline{W}_{ic})\,\epsilon_{ct}\mid \mathbf{X}_{ct}]=0$。误差项里既有不随时间变化的 $\alpha_c$，也有逐年变化的 $u_{ct}$，只要其中任何一部分未被度量且与 climate 相关，这个条件就不成立（第~366 页）。土壤质量、灌溉可得性、人口密度都是这样的因素；Table~3 的平衡性检验给出了直接的证据（见结果节）。

\subsection{面板替代与 identification}

新设定把 climate 变量换成逐年 weather（式 (4)，第~367 页）：
\begin{equation}
y_{ct} = \alpha_c + \gamma_t + \mathbf{X}_{ct}'\boldsymbol{\beta} + \sum_i \theta_i f_i(W_{ict}) + u_{ct},
\label{eq:dg-panel}
\end{equation}
被解释变量改为县级农业利润，$\alpha_c$ 为完整的 \index{county fixed effects}county fixed effects，偏好设定把年份效应 $\gamma_t$ 换成 \index{state-by-year fixed effects}state$\times$year fixed effects $\gamma_{st}$（第~367 页）。被解释变量为何从地价换成利润：``land values capitalize long-run characteristics of sites and, conditional on county fixed effects, annual realizations of weather should not affect land values. Weather does, however, affect farm revenues and expenditures, and their difference is equal to profits''\citep[第~367 页]{deschenes2007economic}。climate 变量为何必须退场：``Second, it is impossible to estimate the effect of the long-run climate averages in a model with county fixed effects, because there is no temporal variation in $\overline{W}_{ic}$. Consequently, we replace the climate variables with annual realizations of weather, $W_{ict}$''\citep[第~367 页]{deschenes2007economic}。

identification 依赖
\begin{equation*}
\E\bigl[f_i(W_{ict})\,u_{ct} \mid \mathbf{X}_{ct}, \alpha_c, \gamma_{st}\bigr] = 0 ,
\end{equation*}
即 weather 参数由 county-specific deviations in weather about the county averages（再调整掉州内各县共同的冲击）识别；作者称这一 variation``is presumed to be orthogonal to unobserved determinants of agricultural profits''\citep[第~357、367 页]{deschenes2007economic}。推断上，全文报告 Eicker--White 异方差稳健标准误，偏好的 fixed effects 设定另报告 Conley 空间相关标准误（经纬度各 7 度、约 500 英里截断；第~366 页及脚注~15）；利润方程以农地总面积的平方根加权（第~374 页）。

\begin{pitfall}[fixed effects 放大 measurement error]
\label{pit:dg-attenuation}
作者自己提醒：``A shortcoming of this approach is that all the fixed effects are likely to magnify the importance of misspecification due to measurement error, which generally attenuates the estimated parameters''\citep[第~367 页]{deschenes2007economic}。fixed effects 削去大部分 variation 之后，weather 变量的 measurement error（测量误差）在剩余 variation 中的占比上升，\index{attenuation bias}attenuation bias（衰减偏误，系数向零偏）随之加重。读 Table~5--6 的系数时必须带着这一条：面板系数偏小，未必是「影响确实小」。
\end{pitfall}

\subsection{理论贡献：这套估计量测到的是什么（式 (1)(2)）}

\S I.B 把「用年际波动推断 climate 效应」的经济学摊开。代表性农户生产给定作物、不能随 weather 换作物，利润为（式 (1)，第~358 页）
\begin{equation}
\pi = p(q(w))\,q(w) - c(q(w)),
\label{eq:dg-profit}
\end{equation}
其中 $q$ 是产出，$w$ 是 weather；价格与总成本是产量的函数，产量是 weather 的函数。对 $w$ 求导（式 (2)，第~358 页）：
\begin{equation}
\frac{\partial \pi}{\partial w} =
\underbrace{\frac{\partial p}{\partial q}\frac{\partial q}{\partial w}\,q}_{\text{价格效应}}
+ \underbrace{\Bigl(p - \frac{\partial c}{\partial q}\Bigr)\frac{\partial q}{\partial w}}_{\text{数量效应}} .
\label{eq:dg-decomp}
\end{equation}
第一项是 weather 冲击经由价格变化带来的利润变化（乘以初始产量 $q$）：减产推高价格，这在短期内缓和代表性农户的损失。第二项是论文想隔离的效应（holding constant crop choice，第~358 页）。关键在两种时间尺度下价格反应的对比：长期里其他农户可以扩产、新农户可以进入，``it is sensible to assume that $(\partial p/\partial q)_{\text{Long Run}} > (\partial p/\partial q)_{\text{Short Run}}$ and is perhaps even equal to zero''\citep[第~358 页]{deschenes2007economic}（$\partial p/\partial q<0$，「$>$」指长期的值更接近零）。所以第一项短期可为正，长期则小得多甚至为零（第~358 页）。state$\times$year fixed effects 的作用之一就是吸收价格的共同变动，使估计接近第二项。

于是得到论文最要紧的一段自我定位：``This paper's empirical strategy relies on year-to-year variation in weather, and thus it is unlikely that farmers are able to switch crops upon a year's weather realization. The import for the subsequent analysis is that our estimates of the impact of climate change may be downward biased, relative to the preferred long-run effect that allows for all economic substitutions. If the degree of climate change is `small', however, our estimates are equal to the preferred long-run effect''\citep[第~359 页]{deschenes2007economic}。Figure~2 的几何版本：作物 1 的利润函数在 $T_1$ 到 $T_2$ 间最优、作物 2 在 $T_2$ 到 $T_3$ 间最优，hedonic 均衡是两段的上包络；climate 从 $T_1$ 升到 $T_3$，能换作物的长期利润读在 $C$ 点，被锁定在作物 1 的短期利润读在 $C'$ 点，$C'-A$ 是 $C-A$ 的向下偏估计（第~359 页）。\textbf{这是 panel approach 与 long-run effect 之间差距的第一次明确陈述}（只给了方向，没有给大小），也是 \chref{ch:merel-gammans-2021} 全部问题的起点。

\begin{intuition}[一年的 weather 与一地的 climate：哪些调整来得及]
D-G 区分的是调整所需的时间。面对 climate 的变化，农户可以调整化肥用量、改变作物组合，甚至把农地改作他用（作者举的例子是住宅开发；第~355 页）；面对某一年的 weather 实现，换作物这类调整来不及做（第~359 页），但投入品的组合仍可调整，作者举的例子是化肥与灌溉用水（第~360 页），所以面板估计仍优于完全不容许 adaptation 的 production function 估计。年际 weather 波动识别的是「作物选择给定、投入可调」之下的响应。\chref{ch:merel-gammans-2021} 用灌渠与风扇的例子表达的是同一区分。
\end{intuition}

\section{结果与解释}

按论文自己的顺序：先检验 hedonic 的前提，再重估 hedonic，最后给出面板估计。

\subsection{平衡性检验：climate normals 与可观测变量不独立}

Table~3（第~368 页）把县按 climate normals 分四分位组，比较各组可观测变量的均值。以 7 月温度为例，每英亩农地价值在四组依次为 1{,}118、1{,}771、1{,}609、1{,}481 美元（$F$ 统计量 25.5），人口密度与人均收入的组间差异同样显著（$F$ 统计量 19.7 与 14.1），均拒绝各组均值相等。作者先说明为什么平衡性值得检验：若可观测变量在各组间平衡，就有两点理由相信 hedonic 估计可靠——``First, consistent inference will not depend on functional form assumptions on the relation between the observable confounders and farm values. Second, the unobservables may be more likely to be balanced''\citep[第~368 页]{deschenes2007economic}。检验结果是不平衡：扩展检验中，112 个解释变量均值相等的原假设有 111 个在 1\% 水平被拒（第~369 页）。作者特别指出人口密度（可视作城市化程度或农地转为住宅的可能性的度量）与农地价值相关，而它对农业单产并无直接影响，这削弱了 hedonic 路线的可信度（第~369 页）。

\subsection{hedonic 重估：不稳健}

复刻 SHF（2005）的设定与数据，基准情景（升温 5\,$^\circ$F、降水 +8\%）对农地价值的总影响为 $-\$543.7$ billion（cropland 权重）或 $+\$69.1$ billion（crop revenue 权重）（第~369 页）。作者再换两组协变量：只保留 climate 变量、去掉其余协变量，估计变为 $-\$98.5$ 与 $+\$437.6$ billion；在 SHF 设定上加 state fixed effects，得到 $-\$477.8$ 与 $+\$1{,}034$ billion——后一个数字``seems implausible, since it is nearly as large as the entire value of agricultural land and buildings in the United States, which was \$1,115 billion in 2002''\citep[第~369 页]{deschenes2007economic}。改用农地面积加权并重构 SHF 数据后，三种设定的估计依次为 $+\$225.1$、$-\$315.4$、$-\$0.6$ billion（第~369 页）。作者结论：``It seems reasonable to conclude that the application of the hedonic approach to the SHF data fails to produce robust estimates of the impact of climate change, even with a single year of data''\citep[第~369 页]{deschenes2007economic}。

D-G 自己的更完整 hedonic（Table~4，第~370--371 页）把同样的不稳定摆得更齐：42 个估计中，偏好的 B、C 两组设定下分年度的点估计从 $-\$202$ 到 $\$321$ billion；六种合并设定与权重组合下从 $-\$75.1$ 到 $\$110.8$ billion；D-G 偏好的 C 设定给出显著为正的 $\$95.2$ 与 $\$110.8$ billion（第~371 页）。作者的评价是：``These results are troubling, because there is no ex ante reason to believe that the estimates from a particular year are more reliable than those from other years''\citep[第~371 页]{deschenes2007economic}。\textbf{注意读法}：这一小节的主要发现是同一数据稍换口径结果就翻转，即不可靠本身（摘要：``to be unreliable because it produces estimates that are extremely sensitive to seemingly minor choices about control variables, sample, and weighting''\citep[第~354 页]{deschenes2007economic}）；作者也顺带指出，其中较可信的 hedonic 设定显示 climate change 对美国农业总体略有益处（第~372 页），但没有把它当作主要结论。

\subsection{论文的 panel approach：利润的影响小且不显著}

Table~5（第~375 页）报告三个情景（基准、Hadley 2 的 2020--2049、Hadley 2 的 2070--2099）下利润的预测变化。偏好设定（第 4 列：土壤控制变量加 state$\times$year fixed effects）在 Hadley 2 长期情景下的结果是``climate change will lead to an increase in agricultural sector profits of roughly \$1.3 billion. This estimate is statistically indistinguishable from zero with either the Eicker-White (parentheses) or Conley (square brackets) standard errors''\citep[第~374 页]{deschenes2007economic}；用较小的标准误，95\% 置信区间约为 $-\$0.5$ 到 $\$3.1$ billion（第~374 页），相对 1987--2002 年均利润 \$32.3 billion 即 $-1.5\%$ 到 $+9.6\%$（第~377 页）。分解上，影响几乎全部来自非灌溉县（第~374 页；Table~5 第 4 列：非灌溉县 1.10、灌溉县 0.18，单位 billion），按变量分则主要来自降水（degree-days 0.39、降水 0.89）。Table~6（第~376 页）共十二行 robustness check，包括线性、三次式等温度与降水的其他函数形式、harmful degree-days、剔除强影响点、灌溉县定义的分界改为 5\% 与 15\%（基准为 10\%）、growing season 的定义与不加权等，估计全部落在 $-\$0.98$ 到 $\$3.59$ billion 之间；对前三行的函数形式变体，作者写明``a zero effect cannot be rejected in any of the cases''\citep[第~376 页]{deschenes2007economic}。分州看（Table~7，第~377--378 页），最大的输家是 California（$-\$750$ million，约其农业利润的 15\%），最大的赢家是 South Dakota（$+\$720$ million）与 Georgia（$+\$540$ million）；分州估计总体不精确，只有八个州能在 5\% 水平拒绝零影响（第~377 页）。单产端（Table~8，第~379 页），偏好设定下玉米与大豆全国产量的预测变化为 $+0.01$ 与 $+0.02$（单位：十亿蒲式耳），都统计上不异于零；作者据此认为利润影响小不太可能只是短期价格上涨所致（第~380 页）。

贯穿本章的例子在这里收尾。数据节 Table~2C 的量级对比说明外推所需的偏离在历史数据中出现过（degree-days 落在尾部，降水宽裕，第~365 页）；引言对 Figure~1 的概括则给出响应面的形状：``The striking finding is that all of the response surfaces are flat over the ranges of the predicted changes in degree-days and inches. If anything, climate change appears to be slightly beneficial for profits and yields''\citep[第~355 页]{deschenes2007economic}。两者合起来就是基准结论的直观版本：在历史上观测得到的偏离范围内，利润与单产对 degree-days 和降水的响应是平的。

\begin{pitfall}[从「影响不显著」读出「没有影响」]
D-G 的结论是``economically and statistically small changes in profits''\citep[第~377 页]{deschenes2007economic}，配一个置信区间（$-\$0.5$ 到 $\$3.1$ billion，第~377 页），不是影响为零的点断言。同样不可忘的是易错点~\ref{pit:dg-attenuation}：fixed effects 放大 measurement error，估计偏向零。后续文献对论文的批评（见「局限与批评」节）恰恰不是「估出了零」，而是「这套 variation 只测到 short-run response」——两件事不要混。
\end{pitfall}

\section{机制与讨论}

福利换算：假定 +\$1.3 billion 的年度利润增量对未来所有年份成立，并取 5\% 的折现率，``This implies that climate change increases the present value of the stream of land rents by about \$26 billion. The 95-percent confidence interval is $-\$10$ billion to \$62 billion''\citep[第~380 页]{deschenes2007economic}。作者强调这个区间``is much tighter than the range of point estimates from the hedonic approach''\citep[第~380 页]{deschenes2007economic}——面板路线的收获是区间的可信性，不只是点估计。换算的两个前提（利润增量对所有未来年份成立、5\% 折现率）是原文正文写明的 assumption，不在表内；引用 \$26 billion 时应一并交代。

\section{局限与批评}

论文自己承认的（\S V，第~380--381 页）。先是四条一般性的 caveats：其一，极端事件（干旱、洪水）与 climate 变率的变化不在 growing season degree-days 与降水的刻画里，论文的分析对这类事件``uninformative''\citep[第~380 页]{deschenes2007economic}；其二，climate change 是全球性的，会改变长期生产成本与相对价格，而年际波动加 state$\times$year fixed effects 的估计不含这类价格变动（作者同时指出 hedonic 同样处理不了）；其三，政府补贴项目影响利润与地价（引 Kirwan 2005）；其四，利润口径低估资本成本（未计自有资本的利息）且未扣无工资报酬的自有劳动。随后是两条会使损害被\textbf{高估}的 caveats：年度波动容不下全部 adaptation（``farmers will undertake these adaptations only if the benefits exceed the costs''\citep[第~381 页]{deschenes2007economic}，所以偏误方向明确），以及 CO$_2$ 浓度上升的增产效应未入模型（第~381 页）。

后续文献指出的：\citet{merel2021climate} 在第~1207 页把这条批评线概括为 panel approach 依赖 weather fluctuations 而非 climatic differences，因而可能反映不了 long-run adaptation；D-G 在 \S I.B 的自认（第~359 页）正是其源头。\citet{hsiang2016climate} 把这套设计的前提形式化为 marginal treatment comparability assumption（第~6 页，见 \chref{ch:hsiang-2016}）；\chref{ch:merel-gammans-2021} 推导出漏掉 penalty term 的面板系数相对 long-run response 的偏差，并给出同时估计 short-run 与 long-run response 的增广设定。

本讲义的评注：D-G 的地位不主要在它的点估计，而在三件事——把 identification 从截面挪到组内、把利润口径的测度缺陷交代到底、以及把「面板测不到全部 long-run 调整」写成了论文自己的结论。读这篇论文要同时拿到这三件；漏掉第三件，就看不出方法线下一步（\chref{ch:burke-hsiang-miguel-2015}、\chref{ch:merel-gammans-2021}）所回应的问题从哪里来。

\section{在本领域脉络中的位置}

论文回答的旧问题：截面 Ricardian 路线的 omitted variables bias 严重到什么程度（答案：hedonic 估计对口径选择极敏感，不可靠）。它开的新问题：panel approach 识别的 short-run response 与 long-run effect 差多少。后来的文献从三个方向处理这个差距：Burke 与 Emerick（2016）的 \index{long differences}long differences（时间线 2016）、Hsiang（2016）对 marginal treatment comparability 的部分检验（\chref{ch:hsiang-2016}）、Mérel 与 Gammans（2021）的 climate penalty 增广回归（\chref{ch:merel-gammans-2021}）。本讲义的归纳：在 \chref{ch:info-climate-change-adaptation} 的方法线上它是 panel 线的第一站，其自认局限（第~359 页）是该线此后各站的出发点。

% 章末参考文献（只列本章引用的条目），勿删
\chapterrefs

% 精读章: Burke, Hsiang & Miguel, Global Non-Linear Effect of Temperature
%         on Economic Production (Nature 2015)
% 原文出处: doi:10.1038/nature15725（所据版本：期刊已刊版 Nature 527: 235--239；
%          本仓 papers/01_methods/2015-Nature-Burke-Hsiang-Miguel-*.pdf 为团队正典，
%          含正文 5 页与 Extended Data Fig. 1--6、Extended Data Table 1--3；
%          本章页码引注均指该版。Supplementary Information（下称 SI）2026-11-19 自
%          作者网站下载，存 papers/01_methods/2015-Nature-...-Supplementary-Information.pdf，
%          该副本数学字体损坏，公式字母不可辨；引 SI 处写节号）
% 体例注记: 2026-10-21 起，正文中的 terms/expressions 一律用英文原词，
%          中文仅保留句法成分；定义处随文给中文对照并登记英文 \index。
% 建档/更新日期由章首 \chdate 渲染（grep: chdate）
% 主笔: 待认领；起草: Claude（2026-11-04，bjr 审订前稿）
% 核对: 2026-11-19 对照本地 PDF 全章重核——原句引用逐条改回原文字面；
%       改正 0.29（原误作 0.28）、75% 所属情景、预测条数、Figure 1 记号、
%       robustness check 清单、损害函数两点差异等处；数据来源据正文参考文献补齐。
% partinfo 定位: nonlinear panel approach 的代表性应用——fixed effects 加二次
%               温度项使横截面 weather 变异进入 identification，见
%               \chref{ch:info-climate-change-adaptation} 的方法线与速查表。

\chapter{Global Non-Linear Effect of Temperature on Economic Production（Burke, Hsiang \& Miguel，Nature 2015）}
\label{ch:burke-hsiang-miguel-2015}
% 章题过长，页眉用短题（preamble 的 \chaptermark 经 \markboth 排页眉）
\chaptermark{Global Non-Linear Temperature Effects（Burke et al.，2015）}
\chdate{2026-11-04}{2026-11-19}

\section{这篇论文在说什么}

一句话概括：论文用 166 国 1960--2010 的年度面板，以``comparing a country to itself in years when it is exposed to warmer- versus cooler-than-average temperatures''\citep[第~236 页]{burke2015global}的组内比较配合二次温度项，发现``country-level economic production is smooth, non-linear, and concave in temperature (Fig.~2a), with a maximum at 13~$^\circ$C''\citep[第~236 页]{burke2015global}，并报告样本期内的变暖没有改变这一响应——``substantial observed warming over the period apparently did not induce notable adaptation''\citep[第~237 页]{burke2015global}，这一判断由此进入领域。

它在 partinfo 里的坐标：nonlinear panel approach 的代表应用站（\chref{ch:info-climate-change-adaptation} 方法线 panel 线第二站）。它被 \citet{merel2021climate} 列入承认 nonlinear panel proposition 的研究名单（第~1208 页），其面板数据也被 Mérel 与 Gammans 用作经验应用素材（第~1224 页），见 \chref{ch:merel-gammans-2021}。

骨架表（Nature Letter 无编号小节；页码为期刊页码；Extended Data 排在正文之后，无期刊页码）：

\begin{table}[htbp]
\centering\small
\begin{tabular}{@{}p{3.4cm}p{1.4cm}p{8.0cm}@{}}
\toprule
段落功能 & 页码 & 讲什么 \\
\midrule
摘要 & 235 & 非线性、13~$^\circ$C 最优、全球通用；projection：不减缓的变暖使 2100 年全球平均收入降约 23\% \\
问题与理论框架 & 235--236 & micro 高度非线性响应如何加总成 macro 平滑凹响应；式 (1)；两条预测 \\
数据与识别 & 236 & 166 国 1960--2010；一国自身冷暖年份的比较；增长的分解 (1)--(4) \\
主要结果 & 236--237 & 13~$^\circ$C 最优、robustness check、滞后效应、rich/poor、农业与非农业 \\
adaptation 的证据 & 237 & 1960--1989 对 1990--2010 两段子样本；未见明显 adaptation \\
Projection & 237--239 & RCP8.5 与 SSP 情景、收入分布、与 IAM 损害函数对比 \\
Extended Data & 无页码 & Fig.~1--6 与 Table~1--3：各项 robustness check 的图表 \\
\bottomrule
\end{tabular}
\end{table}

\section{预备知识}

读者需备两件。其一，增长率回归的用途：论文估计的是``how economic production changes relative to the previous year—that is, annual economic growth—to purge the data of secular factors in each economy that evolve gradually''\citep[第~236 页]{burke2015global}，即以相对上年的变化为对象，把缓慢演变的国别因素差分掉，再用 fixed effects 与趋势项吸收其余。其二，加总记号（第~236 页）：$f_i(T)$ 是行业 $i$ 中单个生产单元（如一个厂商）的生产率对瞬时（如逐日）温度 $T$ 的响应；对给定的国家、时期与行业，$g_i(T-\bar{T})$ 是全部 unit-hours 在各温度上的分布，以平均温度 $\bar{T}$ 为中心；$m_{i1}$ 与 $m_{i2}$ 分别是低于与高于临界温度阈值的 unit-hours 份额。

\section{数据与测度}

产出数据覆盖 166 国、1960--2010 年（第~236 页），取自 World Bank 的 World Development Indicators 2012（正文参考文献~17）；被解释变量是人均 GDP 对数的年度变化，即增长率（Figure~2a 图注，第~236 页；$N=6{,}584$）。climate 变量是年平均温度与降水，取自 Matsuura 与 Willmott 的 Terrestrial Air Temperature and Precipitation 月度与年度序列 v.~3.01（正文参考文献~18）；Extended Data Table~1 表注写明温度以 $^\circ$C、降水以米计。聚合方式写在原文 SI 的 B.1 节：原始数据是 0.5 度格点的月平均温度与月降水总量，按 2000 年的人口密度加权聚合到国家—年份（权重数据取自 Gridded Population of the World）；robustness check 所用的另一套收入数据是 Penn World Tables 8.0；全样本共 6{,}584 个国家—年份观测。

识别所需的全部控制写在一组分解里：``We deconvolve economic growth to account for: (1) all constant differences between countries, for example, culture or history; (2) all common contemporaneous shocks, for example, global price changes or technological innovations; (3) country-specific quadratic trends in growth rates, which may arise, for example, from changing political institutions or economic policies; and (4) the possibly non-linear effects of annual average temperature and rainfall''\citep[第~236 页]{burke2015global}。作者的选型理由是：``This approach is more reliable than only adjusting for observed variables because it accounts for unobserved time-invariant and time-trending covariates''\citep[第~236 页]{burke2015global}。rich 与 poor 的划分用于分组响应：按中位数分为两组（Figure~2b 图注，第~236 页）。SI 的 C.3 节写明分组依据：一国 1980 年经购买力平价（PPP）调整的人均收入低于全球中位数即归入 poor，取 1980 年是因为多数样本国家从这一年起才有 PPP 数据。

\section{模型设定与识别策略：micro 响应如何加总出 macro 响应}

本节是论文的理论贡献。micro 层面的证据早已表明：劳动供给、劳动生产率、作物单产等对温度``exhibit highly non-linear responses''\citep[第~235 页]{burke2015global}，在 20--30\,$^\circ$C 之间存在阈值。而既有的宏观研究报告的是另一幅图景：富国的总体生产率对温度没有响应，穷国只有线性响应（第~235 页）。论文把疑问表述为：``However, it is unclear how these abrupt declines at the micro level are reflected in coarser macro-level data''\citep[第~235--236 页]{burke2015global}。它的答案是加总：若单个 unit-hour 的生产率损失对其他单元影响有限，一国的总产出就是各行业产出之和，每个行业的产出是对该国该时期全部 unit-hours 的积分（式 (1)，第~236 页）
\begin{equation}
Y(\bar{T}) = \sum_i Y_i(\bar{T}) = \sum_i \int_{-\infty}^{\infty} f_i(T)\, g_i(T-\bar{T})\,\mathrm{d}T .
\label{eq:bhm-agg}
\end{equation}
平均温度 $\bar{T}$ 上升时，高于阈值的份额 $m_{i2}$ 逐渐增大，总产出的损失是渐进的、逐步加重的。式~\ref{eq:bhm-agg} 给出两条预测（第~236 页）。其一，$Y(\bar T)$ 是 $\bar{T}$ 的光滑凹函数，其导数等于各 $f_i(T)$ 的导数按各行业在各日温度上的 unit-hours 数加权的平均；micro 响应在阈值处的折点（原文称 kinked）经分布 $g_i$ 加权后被抹平。其二，$Y(\bar T)$ 的峰值所在温度低于 $f_i(T)$ 的阈值，但这是有条件的：``if the slope of $f_i(T)$ above the threshold is steeper than minus the slope below the threshold''\citep[第~236 页]{burke2015global}，作者称 micro 证据支持这一条件。Figure~1（第~235 页）逐步画了这一映射；其 f 图把加总后的斜率标作 $m_{i1}b_{i1}+m_{i2}b_{i2}$，即阈值两侧 micro 斜率按 unit-hours 份额的加权和。

identification 与 \chref{ch:deschenes-greenstone-2007} 同族，但换了尺度：``In practice, we can approximate this experiment by comparing a country to itself in years when it is exposed to warmer- versus cooler-than-average temperatures due to naturally occurring stochastic atmospheric changes''\citep[第~236 页]{burke2015global}。作者明确拒绝截面比较：``We do not compare output across different countries because such comparisons are probably confounded, distinguishing our approach from cross-sectional studies that attribute differences across countries to their temperatures''\citep[第~236 页]{burke2015global}。响应函数取温度的二次式（Extended Data Table~1 的 Temp. 与 Temp.\ sq.\ 两行；基准列的系数为 0.0127 与 $-0.0005$，对应最优温度 13.06\,$^\circ$C）；估计方程是 SI 的 Supplementary Equation (15)（B.2 节）：
\begin{equation}
\Delta Y_{it} = h(T_{it}) + \lambda_1 P_{it} + \lambda_2 P_{it}^2 + \mu_i + \nu_t + \theta_i t + \theta_{i2} t^2 + \eps_{it},
\label{eq:bhm-reg}
\end{equation}
其中 $i$ 为国家、$t$ 为年份，$\Delta Y_{it}$ 是实际人均 GDP 自然对数的一阶差分，$T_{it}$ 与 $P_{it}$ 是年平均温度与年降水，$\mu_i$ 与 $\nu_t$ 是 country 与 year fixed effects，$\theta_i t+\theta_{i2}t^2$ 是各国自己的二次时间趋势；基准设定中 $h(T_{it})$ 取 $T_{it}$ 的二次式。各项的结构与下标依 SI 原式；希腊字母为本讲义所选，因为本仓所收 SI 副本（作者网站版）的数学字体损坏，字母无法辨认。

\begin{ident}[BHM 增长率回归的可搬用设定]
估计方程即式~\ref{eq:bhm-reg}。被解释变量：人均 GDP 对数的年度变化；回归元：年平均温度及其平方、降水及其平方；控制：country fixed effects、year fixed effects、国家二次时间趋势；标准误在国家层面聚类（以上亦见 Extended Data Table~1 表注）。作者对所用 variation 的概括是：考察各国增长率对本国趋势的偏离，是否与本国温度、降水对各自趋势的偏离非线性相关，并扣除各国共同的冲击（第~236 页）。identification assumption 论文没有写成公式；按这一设定，它是「给定上述控制，一国某年的温度偏离与影响增长的其他因素不相关」（本讲义的归纳）。口径要求：温度为全年平均温度，不是 growing season 的温度；产出取增长率而非水平值。
\end{ident}

\begin{intuition}[为什么加总会压低最优温度]
单个生产单元的响应在阈值处由平转陡。一个国家的 unit-hours 沿温度分布摊开，平均温度每升一点，就有一部分 unit-hours 越过阈值，落到损失斜率更陡的一侧。只要阈值以上的斜率比阈值以下的斜率（取负号后）更陡，总产出就会在平均温度尚未达到 micro 阈值时见顶。所以国家层面的最优温度低于 micro 阈值；论文报告的 13~$^\circ$C 与这一预测一致（第~236 页），但式~\ref{eq:bhm-agg} 本身并不给出 13 这个具体数值。
\end{intuition}

\section{结果与解释}

\textbf{主结果。}``We find country-level economic production is smooth, non-linear, and concave in temperature (Fig.~2a), with a maximum at 13~$^\circ$C, well below the threshold values recovered in micro-level analyses and consistent with predictions from equation (1)''\citep[第~236 页]{burke2015global}。响应的形态：``Cold-country productivity increases as annual temperature increases, until the optimum. Productivity declines gradually with further warming, and this decline accelerates at higher temperatures''\citep[第~236 页]{burke2015global}。

\textbf{Robustness check。}作者列出的检验有（第~236--237 页，图表见 Extended Data Fig.~1 与 Extended Data Table~1）：结果具有全球代表性、不由异常值驱动；允许各国的响应随其变富而改变（温度与平均收入交互）；改用高阶多项式或 restricted cubic splines 刻画温度效应；剔除观测少的国家；剔除主要产油国；剔除中国与美国；加入各大洲逐年的共同冲击（continent$\times$year fixed effects）；放松对增长趋势的假定；加入增长率的多期滞后；改用另一套经济数据（Penn World Tables）。温度冲击的滞后效应也考察过，结论是损失之后没有回补：``This `no catch up' behaviour is consistent with the observed response to other climatological disturbances, such as tropical cyclones''\citep[第~237 页]{burke2015global}。

\textbf{rich/poor 与部门分解。}穷国的响应更大——``Poor tropical countries exhibit larger responses mainly because they are hotter on average, not because they are poorer''\citep[第~237 页]{burke2015global}。有迹象显示富国受温度的影响可能略小，但``their response is statistically indistinguishable from poor countries at all temperatures''\citep[第~237 页]{burke2015global}。富国中高温国家少，高温的总效应估计得不精确；即便如此，作者估计有``an 80\% likelihood that the marginal effect of warming is negative at high temperatures in these countries''\citep[第~237 页]{burke2015global}。农业与非农业的总产出对年平均温度都呈非线性，rich 与 poor 两组皆然（第~237 页；Figure~2d、e 在第~236 页）；作者同时说明，低温端在这些子样本里没有显著效应，高温端的负效应在穷国各情形下显著，在富国显著或边际显著（第~237 页）。

\textbf{adaptation 的经验判断。}这是论文进入方法线争论的关键证据：把样本切成 1960--1989 与 1990--2010 两段，``Results using data from 1960–1989 and 1990–2010 are nearly identical (Fig.~2c)''\citep[第~237 页]{burke2015global}；作者的表述是``We do not find that technological advances or the accumulation of wealth and experience since 1960 has fundamentally altered the relationship between productivity and temperature''\citep[第~237 页]{burke2015global}，并进一步说``We find only weak suggestive evidence that richer populations are less vulnerable to warming, and no evidence that experience with high temperatures or technological advances since 1960 have altered the global response to temperature. This suggests that adaptation to climatic change may be more difficult than previously believed''\citep[第~237 页]{burke2015global}。\textbf{注意这条证据的口径}：它是两段子样本之间的系数对比，不是 adaptation 的直接测度；它检验的是「全球响应函数在两段时期之间是否移动」。而按 \chref{ch:merel-gammans-2021} 的分析，这条由二次面板估计出的响应函数本身已经是 short-run 与 long-run response 的加权平均。读这一结论时两层限定要叠加考虑。

\textbf{Projection。}把估计的非线性响应代入 RCP8.5 的变暖情景与 SSP 的基线增长情景（第~237 页）。各国的升温幅度取 CMIP5 全部模式在 RCP8.5 下的集合平均，按人口加权聚合到国家，样本国家到 2100 年的升温在 2.7 到 5.8\,$^\circ$C 之间（SI 的 D.1 节）。基准模型下，``climate change reduces projected global output by 23\% in 2100 (best estimate, SSP5) relative to a world without climate change, although statistical uncertainty allows for positive impacts with probability 0.29''\citep[第~238 页]{burke2015global}。收入分布的后果更重，同样在基准估计下，``average income in the poorest 40\% of countries declines 75\% by 2100 relative to a world without climate change, while the richest 20\% experience slight gains''\citep[第~238 页]{burke2015global}。projection 的前提由作者自己交代：``This approach assumes future economies respond to temperature changes similarly to today's economies—perhaps a reasonable assumption given the observed lack of adaptation during our 50-year sample''\citep[第~237 页]{burke2015global}。与 IAM 损害函数的对比，摘要的概括是：``expected global losses are approximately linear in global mean temperature, with median losses many times larger than leading models indicate''\citep[第~235 页]{burke2015global}。

\section{机制与讨论}

加总机制（式~\ref{eq:bhm-agg}）在结果端的对应是：macro 最优温度 13~$^\circ$C 远低于 micro 研究得到的阈值（20--30\,$^\circ$C），作者称这与式 (1) 的预测一致（第~235--236 页）。不平等机制的表述：``Strong negative correlation between baseline income and baseline temperature indicates that warming may amplify global inequality because hot, poor countries will probably suffer the largest reduction in growth''\citep[第~238 页]{burke2015global}。对 IAM 的含义：作者用自己的结果构造经验损害函数，与 IAM 中常用的损害函数比较，``Commonly used functions are within our estimated uncertainty, but differ in two important respects''\citep[第~238 页]{burke2015global}。两点差异写在第~239 页：其一，论文预测的全球损失对温度大致呈线性（略凹），不是此前理论设定的二次或指数形式；其二，损害函数的斜率在小幅升温时已经很大，2\,$^\circ$C 升温的预期损失是既有估计的 2.5 到 100 倍，更高升温下至少是 2.5 倍。

\section{局限与批评}

论文自己承认的：其一，projection 假定未来经济体对温度的响应与今天相同，其辩护是样本期内未见 adaptation（第~237 页）；作者在结尾处两面都说了，前所未有的创新或防御性投资可能减轻这些影响，社会冲突或贸易受阻则可能加重（第~239 页）。其二，rich 与 poor 的差异估计得不精确：高温富国少，富国高温效应的不确定性大；作者能说的只是``the non-linearity of the rich-country response alone is statistically significant''\citep[第~237 页]{burke2015global}（$P<0.1$）与上文 80\% 的可能性，区分 rich 与 poor 响应的 projection 也因每组数据更少而更不确定（第~237 页）。其三，估计只含温度效应，不含 climate change 的其他损失来源，如热带气旋与海平面上升（第~239 页）。其四，统计不确定性不小，作者处处给出显式的概率（如基准估计下全球影响为正的概率 0.29，第~238 页）。

后续文献指出的：\citet{merel2021climate} 的偏差分解为解读论文的系数补上了缺失的一环——二次面板的温度系数是 long-run 与 short-run response 的加权平均（无 skewness 时为 convex combination），long-run 的权重取决于总体 weather variation 中来自横截面 climatic differences 的份额（第~1209 页）。这套国家面板的横截面差异大，long-run 成分的占比应当较高，但具体份额须按 \chref{ch:merel-gammans-2021} 的权重公式计算后才可断言（本讲义的引申；Mérel 与 Gammans 对全球面板的计算在其在线附录 G.4，本仓未收，待补）。对 projection 的外部批评（收入弹性、加总口径等）在后续文献中存在，待精读后补。

本讲义的评注：论文与 \chref{ch:deschenes-greenstone-2007} 的对照是方法线上最适合放在一起讲的一组——同一族组内 identification，两者的回归都含 weather 变量的二次项（D-G 第~367 页；Mérel 与 Gammans 把两篇同列于二次设定面板的例子，第~1208 页脚注~3），差别在尺度与被解释变量：D-G 用美国县级农业利润，BHM 用国家人均 GDP 增长率。BHM 自己引以为旁证的是另一项美国县级研究：``Our finding that rich countries respond non-linearly to temperature is consistent with recent county-level results in the United States''\citep[第~237 页；所引为 Deryugina 与 Hsiang 2014 的工作稿，正文参考文献~8]{burke2015global}。

\section{在本领域脉络中的位置}

论文回答的旧问题：温度与产出的 macro 关系是否非线性、全球是否共用一条响应（答案：是，13~$^\circ$C 最优）。它并非 nonlinear panel approach 的首创：penalty 结构与「hybrid estimator」最早见于 \citet{mcintosh2006identifying}，农业端的二次 weather 面板至少可追溯到 \citet{deschenes2007economic}（第~367 页；名单见 \citealp[第~1208 页脚注~3]{merel2021climate}）；论文的地位是全球宏观尺度上的旗舰应用。它开的新问题：这条响应里有多少 long-run adaptation 成分——由 \chref{ch:merel-gammans-2021} 形式化为可计算量；「样本期内未见 adaptation」的判断也把 adaptation 的检验推到台前（\chref{ch:deschenes-greenstone-2007} 的自认局限、\chref{ch:hsiang-2016} 对各时间尺度估计的比较，处理的是同一争论的不同侧面）。本章在 \chref{ch:info-climate-change-adaptation} 的方法线与速查表中均有登记。

% 章末参考文献（只列本章引用的条目），勿删
\chapterrefs

% 精读章: Hsiang, Climate Econometrics (Annual Review of Resource Economics 2016)
% 原文出处: doi:10.1146/annurev-resource-100815-095343（所据版本：NBER Working
%          Paper 22181, April 2016——本仓 papers/01_methods/2016-NBER-w22181-*.pdf
%          为团队正典，本章页码引注均指该 NBER 版的印刷页码；已刊版 ARRE 8: 43--75
%          并存，页码不同，引用页码须注明所据版本。期刊卷期与 DOI 不见于 NBER 版
%          PDF，出处见 bibs/refs.bib 该条的 note）
% 体例注记: 2026-10-21 起，正文中的 terms/expressions 一律用英文原词，
%          中文仅保留句法成分；定义处随文给中文对照并登记英文 \index。
% 建档/更新日期由章首 \chdate 渲染（grep: chdate）
% 主笔: 待认领；起草: Claude（2026-11-04，bjr 审订前稿）
% 核对: 2026-11-19 对照本地 PDF 全章重核——原句引用逐条改回原文字面；
%       改正式 (16)--(20) 的推导次序、式 (24)、玉米应用的频段、骨架表各节
%       页码，以及 marginal treatment comparability 与 unit homogeneity 的关系。
% partinfo 定位: 把「weather 回归能否识别 climate 效应」的争论正式化，整理
%               cross-section／time-series／long differences 三种 research designs
%               与两件工具（marginal treatment comparability、
%               frequency-identification tradeoff），见
%               \chref{ch:info-climate-change-adaptation} 的方法线与速查表。

\chapter{Climate Econometrics（Hsiang，ARRE 2016）}
\label{ch:hsiang-2016}
\chdate{2026-11-04}{2026-11-19}

\section{这篇论文在说什么}

一句话概括：这篇综述（\citealp{hsiang2016climate}，所据 NBER w22181 版）把``regressions on `weather' variables''\citep[第~2 页]{hsiang2016climate}能否识别 climate 效应这个``particularly thorny''\citep[第~2 页]{hsiang2016climate}的争论正式化，在可以写明的条件下得出``the use of weather variables in regressions is justified and, perhaps surprisingly, dominates traditionally preferred methods''\citep[第~2 页]{hsiang2016climate}；同时整理了 climate econometrics 的三种 research designs 与两件工具（frequency-identification tradeoff、对 marginal treatment comparability 的部分检验），并给出 weather variation 恰好识别 climate 效应的充分条件（Envelope Theorem 与 Gradient Theorem，\S 2.4）。

它在 partinfo 里的坐标：方法线上联结两条线的一站（\chref{ch:info-climate-change-adaptation}）——它把 cross-sectional line 与 panel line 的设计放进同一组记号，逐一列出各自依赖的 assumption。本讲义的归纳：\chref{ch:deschenes-greenstone-2007} 自认的局限与 \chref{ch:burke-hsiang-miguel-2015} 的 adaptation 判断，都可以在这套框架里重新表述；\chref{ch:merel-gammans-2021} 随后自陈与其关系：``our result is both different and more specific''\citep[第~1217 页]{merel2021climate}。

骨架表（所据版本：NBER Working Paper 22181；页码为该版印刷页码。\S 4 及以后本讲义只读到与本 part 直接相关的 \S 4.1.3，其余待补读）：

\begin{table}[htbp]
\centering\small
\begin{tabular}{@{}p{4.4cm}p{1.4cm}p{7.2cm}@{}}
\toprule
原文章节 & 页码 & 讲什么 \\
\midrule
\S 1 Introduction & 2--4 & 1.1 climate 的形式化定义（$\mathbf{C}_{i\tau}$ 对 $\mathbf{c}_{i\tau}$）；1.2 direct 与 belief 两条作用通道 \\
\S 2 The Empirical Problem & 4--14 & 2.1 三种 research designs；2.2 低频 variation 与可信识别的权衡；2.3 marginal treatment comparability 的部分检验；2.4 用 weather variation 恰好识别 climate 效应 \\
\S 3 Measurement of Climate Variables & 15--17 & 从社会与经济的角度压缩 climate 变量的维度；测度创新（Figure~4） \\
\S 4 Econometric Models & 17--25 & 4.1 非线性效应；4.2 位移与滞后；4.3 统计不确定性；4.4 adaptation 的测度；4.5 跨研究的比较与综合 \\
\S 5 Attribution and Projection & 25--26 & 历史归因；未来 climate change 效应的 projection \\
\S 6 Remaining Challenges & 26--27 & 效应与机制的对应、adaptation 与一般均衡、无先例的事件、与理论和数值模型的结合 \\
\S 7 Conclusion & 27--28 & 结论 \\
\bottomrule
\end{tabular}
\end{table}

\section{预备知识}

三件工具，首用交代来历。其一，因果推断的 Fundamental Problem of Causal Inference（Holland 1986，因果推断文献的术语）：同一总体不可能在同一时段同时经历 $\mathbf{C}$ 与 $\mathbf{C}+\Delta\mathbf{C}$，故式~\ref{eq:hs-ate} 的处理效应不可直接观测（\S 2 开头，第~4 页）。其二，\index{envelope theorem}Envelope Theorem（包络定理，微观经济学的标准结果）：最优值函数对参数的导数，等于目标函数在最优点处对该参数的偏导数；选择变量随参数的调整对最优值没有一阶影响。其三，\index{gradient theorem}Gradient Theorem（梯度定理，多元微积分的标准结果）：梯度沿一条路径的线积分等于路径两端函数值之差。后两者在 \S 2.4 使用（第~10--14 页；论文在第~11 页点出这两个定理的名字，但不作陈述，以上表述是教科书的通行说法）。

\section{数据与测度}

综述不依赖单一数据集，本节整理它对「climate 变量怎么测」的要点（\S 3，第~15--17 页）。\S 3 的主题是如何从社会与经济的角度压缩 $\mathbf{c}$ 的维度。早期研究多用粗测度（纬度、事件虚拟变量），作者列举其缺陷，其中之一是``coarse climate measures introduce large measurement errors that will cause attenuation bias, leading to under-rejection of the null hypothesis''\citep[第~15 页]{hsiang2016climate}；另一条是研究者常常关心 \index{dose-response function}dose-response function（剂量—反应函数）$\E[Y\mid\mathbf{c}] = f(\mathbf{c})$ 的结构（式 (25)，第~15 页），粗测度给不出。测度创新的方向见其 Figure~4（第~15 页图注）：用逐日最低、最高温度之间插值出的逐小时温度构造 degree-days（Schlenker 与 Roberts 2009）、用风场模型重建热带气旋暴露（Hsiang 与 Jina 2014）、识别与 ENSO 强耦合的 teleconnected 像素（Hsiang 等 2011）、以下雨天数刻画年内降水的分布（Fishman 2016）。这一节与本讲义对读者的要求一致：climate 变量的口径是跨研究不可比的首要来源。

\section{模型设定与识别策略：weather 回归何时识别 climate 效应}

本节是论文的理论贡献，按原文四步走：定义、问题、三种设计的取舍、两件工具与一个正式结果。

\subsection{第一步：climate 与 weather 的形式化定义}

位置 $i$ 在每个时刻 $t$ 有一个状态向量 $\mathbf{v}_{it} = [$temperature, precipitation, humidity,\dots$]$（式 (1)，第~2 页）；在时段 $\tau$ 内，它从分布 $\psi(\mathbf{C}_{i\tau})$ 中抽出（式 (2)）。刻画该分布的参数向量 $\mathbf{C}_{i\tau}$ 定义为 \index{climate}climate；时段内实际实现的状态有一个经验分布 $\psi(\mathbf{c}_{i\tau})$，刻画它的向量 $\mathbf{c}_{i\tau}$ 定义为 \index{weather}weather（第~2--3 页）——``Define $\mathbf{C}_{i\tau}$ to be the \emph{climate} at $i$ during $\tau$, since it characterizes the distribution of possible realized states $\mathbf{v}_{it}$''\citep[第~2 页]{hsiang2016climate}，``we define $\mathbf{c}_{i\tau}$ to be a description of the \emph{weather} during $\tau$''\citep[第~3 页]{hsiang2016climate}。本 part 的 weather／climate 区分（\chref{ch:info-climate-change-adaptation}）在计量上对应的就是这一对记号。作者还强调 climate 没有固有的时间尺度（通行的 30 年平均定义``is fairly arbitrary''\citep[第~3 页]{hsiang2016climate}）。

climate 影响结果有两条通道：\index{direct effect}direct effect 与 \index{belief effect}belief effect。``the climate during $\tau$ influences what realizations of weather $\mathbf{c}$ actually occur during that interval, which in turn affects a population directly (e.g. a rainy climate generates rain, causing people to get wet)—call this the `direct effect' of climate. Second, individuals' beliefs over the structure of $\mathbf{C}$ may affect their decisions and resulting outcomes, regardless of what $\mathbf{c}$ is realized (e.g. if people believe their climate is rainy, some will buy umbrellas)—refer to this as the `belief effect'''\citep[第~3 页]{hsiang2016climate}。总效应写成 $Y(\mathbf{C}) = Y(\mathbf{c}(\mathbf{C}), \mathbf{b}(\mathbf{C}))$（式 (3)，第~3 页），其梯度按 direct 与 belief 两块分解（式 (4)，第~4 页）；作者注：``interactions between beliefs and direct impacts ($\frac{\partial^2 Y}{\partial b_n\partial c_k}$) and belief effects themselves are together often referred to as `adaptations' in the literature''\citep[第~4 页]{hsiang2016climate}——在这个框架里，adaptation 就是 belief effect 加上它与 direct effect 的交互。

\subsection{第二步：实证问题与两条 assumption}

目标是 climate 变化的平均处理效应（式 (5)，第~4 页）
\begin{equation}
\beta = \E[Y_{i\tau}\mid\mathbf{C}_{i\tau}+\Delta\mathbf{C}_{i\tau}, \mathbf{x}_{i\tau}] - \E[Y_{i\tau}\mid\mathbf{C}_{i\tau}, \mathbf{x}_{i\tau}] ,
\label{eq:hs-ate}
\end{equation}
其中 $\mathbf{x}$ 是可观测的非 climate 因素。同一单元不能同时处于两种 climate 之下，只能拿另一个单元或另一个时段来近似。跨单元比较要求两单元同质：$\E[Y_{i\tau}\mid\mathbf{C},\mathbf{x}_{i\tau}] = \E[Y_{j\tau}\mid\mathbf{C},\mathbf{x}_{j\tau}]$，即 \index{unit homogeneity}unit homogeneity assumption（单位同质假设，式 (6)，第~4--5 页）。时间序列路线只需要它的一个较弱的形式（同一单元跨时期可比），但必须另加第二条 assumption：\index{marginal treatment comparability}marginal treatment comparability assumption（边际处理可比性，第~6 页）——某一时段实现的 weather 相对基准 climate $\mathbf{C}_1$ 的偏离 $\mathbf{c}_{\tau}-\mathbf{C}_1$ 所带来的结果变化，等于 climate 参数本身作同样大小的调整 $\Delta\mathbf{C} = \mathbf{C}_2-\mathbf{C}_1$ 所带来的结果变化（式 (9)，第~6 页）。作者直言这条 assumption``has been widely debated''\citep[第~6 页]{hsiang2016climate}，故后文给检验与条件。

\subsection{第三步：三种 research designs 的取舍}

``There are essentially three research designs in use that approximate the average treatment effect in Eq.~5: cross-sectional approaches, use of time-series variation, and a hybrid known as `long differences'''\citep[第~5 页]{hsiang2016climate}。各自的 variation 与代价：

\textbf{\index{cross-sectional approach}Cross-sectional approach。}同一时期跨单元比较（式 (8)）。长处是隐含地包括了 adaptation：``since farmers who inhabit a location for a long period will have a strong grasp of $\mathbf{C}$ at their location and will adjust farm investments and management to optimize based on these beliefs, farm prices can be assumed to reflect all direct effects as well as all belief effects''\citep[第~5 页；此处在原文中即对 Mendelsohn 等 1994 农地价值路线的概括]{hsiang2016climate}。短处：``A weakness of the cross-sectional approach is its vulnerability to omitted variables bias''\citep[第~5 页]{hsiang2016climate}，且``there exists no systematic method for determining whether any key variables are omitted from Eq.~8 and thus an econometrician can never be certain their model is unbiased''\citep[第~5 页]{hsiang2016climate}。

\textbf{Time-series（含 panel）approach。}单元自身跨期比较（式 (10)，unit fixed effects 加趋势项）。这一思路由 Huntington（1922）提出，他还用工厂工人的数据实施过类似的比较（第~6--7 页）；它获得现代地位是在 Deschênes 与 Greenstone（2007）手中：``This approach gained prominence in modern economic analysis when used by Deschênes \& Greenstone (2007), who analyzed whether agricultural profits in US counties responded to `random fluctuations in weather'''\citep[第~7 页]{hsiang2016climate}。长处：``The core benefit of this approach is that it accounts for unobservable differences between units, eliminating a potential source of omitted variables bias''\citep[第~7 页]{hsiang2016climate}。短处：随时间变化、又与 $\mathbf{c}$ 相关的因素若被遗漏，估计仍会有偏；作者引 Hsiang（2010）的例子——温度、降水与飓风都受 ENSO 调制，漏掉其中重要的 climate 变量，$\hat\beta_{TS}$ 就可能有偏（第~7 页）。

\textbf{\index{long differences}Long differences。}跨单元比较各自在长时段上的变化（式 (11)，第~8 页），定位是折中：``An approach that aims to compromise between the strengths and weaknesses of cross-sectional analysis and time series identification is the long differences strategy, where changes for both the outcome and the climate within locations are correlated across locations''\citep[第~7 页]{hsiang2016climate}；在增长、作物单产与冲突的应用中，研究者``have found that $\beta_{LD}$ is almost identical to $\beta_{TS}$, leading them to conclude that gradual changes in $\mathbf{c}$ likely induce similar effects to more rapid changes in $\mathbf{c}$''\citep[第~8 页，指 Dell et al. 2012、Burke \& Emerick、Lobell \& Asner 2003、Burke et al. 2015b 等的应用]{hsiang2016climate}。

\subsection{第四步：两件工具与一个正式结果}

\textbf{工具一：\index{frequency-identification tradeoff}frequency-identification tradeoff。}belief effect 只对持续的环境变化起作用——只持续五分钟的变化不会让人更新对 climate 的信念（第~8 页）——所以``it is widely thought that low-frequency data ... is required to measure belief effects when using time-series variation''\citep[第~8 页]{hsiang2016climate}；但所用 variation 的频率越低，unit homogeneity 越难成立（相隔数世纪的同一总体未必可比），``it comes at the cost of less credible identification, an issue highlighted by Hsiang \& Burke (2014) as the \emph{frequency-identification tradeoff}''\citep[第~8--9 页]{hsiang2016climate}。缓解条件：若 belief effects 主要通过改变 direct effects 起作用，``total effects of climate still may be nearly identified with high-frequency time-series''\citep[第~9 页]{hsiang2016climate}。

\textbf{工具二：部分检验（\S 2.3）。}若三种设计的估计相等，$\hat\beta_{CS}=\hat\beta_{LD}=\hat\beta_{TS}$，则``one possible explanation is that the marginal treatment comparability assumption is valid and temporary changes in realizations of $\mathbf{c}$ have similar effects to analogous changes in $\mathbf{C}$''\citep[第~9 页]{hsiang2016climate}。它的实施见下文「结果与解释」一节。

\textbf{正式结果：exact identification（\S 2.4）。}这是本综述最重的理论贡献。设结果是最大化问题的值函数（式 (14)，第~11 页）
\begin{equation*}
Y(\mathbf{C}) = \max_{\mathbf{b}\in\mathbb{R}^N} z\bigl(\mathbf{b}, \mathbf{c}(\mathbf{C})\bigr),
\end{equation*}
$z$ 对 $\mathbf{b}$ 可微且为凹。推导分五步（式 (15)--(20)，第~11、13 页）。先对 climate 的第 $k$ 个分量用链式法则，$\mathrm{d}Y/\mathrm{d}C_k$ 拆成三项：$z$ 对 $C_k$ 的直接偏导、经由行动 $b_n$ 的间接项、经由 weather $c_\kappa$ 的间接项（式 (15)）。第一项为零：climate 作为一个概率分布，只能经由它引起的 weather 实现和基于信念的行动影响结果，$\partial z/\partial C_k=0$（式 (16)）。第二项为零：$Y$ 是 $z$ 对全部 adaptation 行动优化后的结果，最优处的一阶条件给出 $\partial z/\partial b_n=0$（式 (17)）。代入后只剩经由 weather 的一项（式 (18)，作者称之为 Envelope Theorem 的应用）。最后只看彼此对应的 weather 与 climate 维度，$\mathrm{d}c_\kappa/\mathrm{d}C_k$ 在 $\kappa=k$ 时为 1、否则为 0（式 (19)），得到式 (20)：
\begin{equation*}
\frac{\mathrm{d}Y(\mathbf{C}_1)}{\mathrm{d}\mathbf{C}_k} = \frac{\partial Y(\mathbf{C}_1)}{\partial \mathbf{c}_k} ,
\end{equation*}
即``Locally, the marginal effect of the climate on $Y$ is identical to the marginal effect of the weather. Eq.~20 implies that Eq.~9, the marginal treatment comparability assumption, holds''\citep[式 (20) 及引文，第~13 页]{hsiang2016climate}。再用 Gradient Theorem 把 marginal effect（边际效应）沿 climate 路径积分，得非边际 climate 变化的效应（式 (23)，第~14 页）
\begin{equation}
Y(\mathbf{C}_2) - Y(\mathbf{C}_1) = \int_{\mathbf{C}_1}^{\mathbf{C}_2} \hat\beta_{TS}\big|_{\mathbf{C}} \cdot \mathrm{d}\mathbf{C} ,
\label{eq:hs-integral}
\end{equation}
``in the presence of adaptation adjustments in $\mathbf{b}$, using only time-series estimates''\citep[式 (23) 及引文，第~14 页]{hsiang2016climate}。作者的总结：``If the outcome is a solution to a maximization problem (Eq.~14) for a function $z(.)$ that is continuous and differentiable in the space of all adaptive actions $\mathbf{b}$, then by application of the Envelope Theorem (Eq.~18) we know that the marginal effect of the climate is exactly the same as the marginal effect of an equally structured change in the weather distribution (Eq.~20), if both are evaluated locally relative to an initial climate. By the Gradient Theorem we know that a sequence of marginal effects of the weather empirically estimated via time-series variation at sequential values of $\mathbf{C}$ can then be integrated to compute the effect of non-marginal climate changes''\citep[第~14 页]{hsiang2016climate}。

这一结果的直觉是论文自带的相邻两户例子（第~11 页；几何版本是其 Figure~3，第~12 页）：一条南北向街道上的两户邻居，其他方面完全相同，只是纬度略异、climate 略异，各自都已针对自己的 climate 做了优化。两户 adaptation 决策的细微差别带来的净效应为零，因为北户多得的边际收益恰被多付的边际成本抵消（两户都处在最大值点）。于是南户若经历北户的 weather，其 weather 效应就等于两户之间的截面差异。\textbf{注意适用范围}：这一等价只在「结果已被优化」的边际处成立；截面路线要恢复 $Y(\mathbf{C})$，还需各单元的积分常数相同（式 (24)，$\mathrm{d}\phi/\mathrm{d}i=0$），``which implies the strong form of the unit homogeneity assumption that units are comparable in levels conditional on the climate (Eq.~6). Thus the set of assumptions necessary for valid cross-sectional identification in this setting is strictly larger than the set of assumptions required for valid time-series identification''\citep[第~14 页]{hsiang2016climate}。

\begin{intuition}[Envelope 论证与 M-G 切点的关系]
Envelope 论证说的是：在一个已经优化的 climate 处，weather 的 marginal effect 等于 climate 的 marginal effect。对「一条在 weather 实现上非线性的面板回归」能恢复什么，综述在 \S 4.1.3 另有讨论（第~19--20 页）：若存在一条全局的非线性响应，这样的回归可以把它估出来，但响应面只在各单元 weather 实现的切点附近由时间序列 variation 局部识别，结论在 weather 实现趋于本地均值的极限下才严格成立。\chref{ch:merel-gammans-2021} 的几何把这一限定讲得更具体：在他们的 climate penalty 模型里，各 location 的 short-run response 在自己的 $\mu_i$ 处与 long-run response 相切（该章图~\ref{fig:cca-lrsr}），切点就是 envelope 条件成立的地方；离开切点，漏掉 penalty term 的面板系数是两种 response 的加权平均（自陈：``our result is both different and more specific''\citep[第~1217 页]{merel2021climate}）。两文合读，这一论证的适用边界才完整。
\end{intuition}

\section{结果与解释：部分检验的实施（玉米单产）}

综述自己的经验证据集中在 \S 2.3 的应用（第~10 页；Figure~2 在第~9 页）。数据：Schlenker 与 Roberts（2009）所用的县级玉米单产、温度与降水数据，更新至 2014 年，限于 100 度经线以东、没有缺失观测的 730 个县。做法：对三个变量的序列施 Baxter--King 带通滤波，按 2--5、6--9、10--13、14--17、18--33 年五个周期段分别估计式 (10)（滤波样本为 1962--2002 年），与原始年度数据的估计（$\hat\beta_{TS}$，即 Schlenker 与 Roberts 的做法）、long differences（$\hat\beta_{LD}$，Burke 与 Emerick）及横截面（$\hat\beta_{CS}$，Schlenker 等 2006）并列。结果：``In all cases, except the cross-section, these estimated effects are near one another and not statistically different, suggesting that variations in temperature over time have similar effects on maize yields in this context, regardless of the time-scale of these variations''\citep[第~10 页]{hsiang2016climate}。横截面例外的原因可辨：``the fact that $\hat\beta_{CS}$ changes substantially (to more closely resemble time-series estimates) when rainfall terms are \emph{omitted} highlights the vulnerability of the cross-sectional approach to misspecification''\citep[第~10 页]{hsiang2016climate}。作者小结：``Nonetheless, these results overall appear consistent with an assumption of marginal treatment comparability in this context, at least for time scales shorter than 33 years''\citep[第~10 页]{hsiang2016climate}。\textbf{读图须知}：检验的说服力直接受各设计的 unit homogeneity assumption 是否成立的约束，``It is theoretically possible to obtain $\hat\beta_{CS}=\hat\beta_{LD}=\hat\beta_{TS}$ by chance even if all key assumptions are violated''\citep[第~10 页]{hsiang2016climate}——各估计相等不能证明 assumption 成立，只能排除一部分反驳。

\section{局限与批评}

论文自己承认的：\S 2.4 的结论依赖三个条件，作者指出既有的经验研究并不检查它们——``that $Y$ is the solution to a [constrained] maximization, that adaptations $\mathbf{b}$ take on continuous values, and that the maximand function $z(.)$ is differentiable in $\mathbf{b}$''\citep[第~14 页]{hsiang2016climate}；而且``many prior studies do not properly compute climate effects via Eq.~23''\citep[第~14 页，作者举 Schlenker et al. 2013 与 Houser et al. 2015 为突出的例外，并提到另有数项研究计算正确或以线性外推近似]{hsiang2016climate}；部分检验中各估计相等可以出于巧合（第~10 页）。\S 6 登记的其余方法论挑战（第~26--27 页）本讲义未读，待补写。后续文献：\chref{ch:merel-gammans-2021} 指出边际处等同只是特例——漏掉 penalty term 的二次面板估计到的是 long-run 与 short-run response 的加权平均，与 long-run response 相切的位置由 weather 的 variation 结构决定（该文第~1215--1219 页）。

本讲义的评注：本综述与 \chref{ch:merel-gammans-2021} 是一组必须对着读的文献。Hsiang 回答「何时可以相信 weather 回归识别了 climate 效应」（边际处、被优化的结果），M-G 回答「在一个具体的函数形式下，偏离有多大、偏向哪边」（由 skewness 与 variation 构成决定）。前者的 Envelope 论证对应后者的切点几何；后者的偏差公式是前者没有给出的量化。

\section{在本领域脉络中的位置}

本讲义的归纳：这篇综述上承两站经验研究留下的问题——\chref{ch:deschenes-greenstone-2007} 自认面板估计不含全部 long-run 调整，\chref{ch:burke-hsiang-miguel-2015} 的非线性面板系数该如何解读尚无定论——它把这些问题放进同一框架，给出三种 designs 各自的 assumption 清单、一个部分检验与一个正式结果。它开的新问题：frequency-identification tradeoff 之下，adaptation 的直接测度如何做（原文 \S 4.4，第~23--24 页，待补读）；\chref{ch:merel-gammans-2021} 的偏差分解可视为在二次设定下对「面板系数含多少 long-run 成分」的第一个闭式回答。本章在 \chref{ch:info-climate-change-adaptation} 的方法线与速查表中均有登记。

% 章末参考文献（只列本章引用的条目），勿删
\chapterrefs

% 精读章: Mérel & Gammans, Climate Econometrics: Can the Panel Approach
%         Account for Long-Run Adaptation? (AJAE 2021)
% 原文出处: doi:10.1111/ajae.12200（所据版本：期刊已刊版，AJAE 103(4): 1207--1238；
%          本仓 papers/01_methods/ 下 PDF 为团队正典，本章页码引注均指该版。
%          在线附录 A--H 不在该 PDF 内，凡涉及附录处只转述正文对附录的说法）
% 版本注记: 第二节原稿的五条编号假设依据工作稿第一版（UC Davis ARE,
%          long-run_effects_in_panel_models_v1）起草，曾按已刊版逐条核对；
%          2026-11-06 起按主笔提纲重写为「数据与方程 → Assumption 1 →
%          interpretation → 几何理解」，编号假设只保留 Assumption 1，
%          optimizing agents 行为推导收进卡片；行为推导、公式 (1)(2)、
%          climate penalty 命名与图 1 均以已刊版为准。
% 体例注记: 2026-10-21 起，正文中的 terms/expressions 一律用英文原词
%          （weather、climate、location/locational、climate penalty……），
%          中文仅保留句法成分；\index 随之登记英文词条。
%          本节按「先 identification 再反推 assumption」组织（2026-10-21，
%          见 CONTRIBUTING「精读章·模型设定与识别策略」）。
% 建档/更新日期由章首 \chdate 渲染（grep: chdate）
% 主笔: bjr；初稿: 2026-10-20
% 续写: 2026-11-04 导读地图、预备知识、渐近偏差分解、一致估计、数据与测度、
%       结果与解释、局限与批评、脉络位置各节（Claude 协助起草，bjr 审订）
% 核对与续写: 2026-11-19 对照本地 PDF 全章重核——原句引用逐条改回原文字面；
%       矩记号、θ 与 μ 的上标、式号、脚注号、Table 1--4 的数字与显著性、
%       反事实三条 remarks 的次序改从原文；按主笔 2026-11-06 批注，
%       「渐近偏差分解」一节的两个小节都改为由 omitted variables bias 公式推出
%       （正文写推导逻辑，卡片写化简过程）。
% partinfo 定位: 把「面板法能否反映长期适应」形式化为可计算的偏差分解，
%               见 \chref{ch:info-climate-change-adaptation} 的方法线与速查表。

\chapter{Climate Econometrics: Can the Panel Approach Account for Long-Run Adaptation?（Mérel \& Gammans，AJAE 2021）}
\label{ch:merel-gammans-2021}
% 章题过长，页眉用短题（preamble 的 \chaptermark 经 \markboth 排页眉）
\chaptermark{Climate Econometrics（Mérel \& Gammans，2021）}
\chdate{2026-10-19}{2026-11-19}

\section{这篇论文在说什么}

一句话概括：论文把「nonlinear panel model 的估计里含有多少 long-run 成分」形式化为可计算的 \index{asymptotic bias}asymptotic bias（渐近偏差）分解。如摘要中所说：``Building upon a framework proposed by McIntosh and Schlenker (2006), this paper identifies empirical conditions under which the nonlinear panel approach can approximate a long-run response to climate. When these conditions fail, the obtained relationship may still be interpretable as a weighted average of underlying short-run and long-run responses''\citep[第~1207 页]{merel2021climate}。论文同时给出一致估计两条 response 的增广回归（climate penalty 设定），并在美国与法国的作物单产面板上量化——温度系数已主要反映 long-run response（长期权重 98\%、88\%、86\%），降水系数含大量 short-run 成分（长期权重 41\%、7\%、53\%）（第~1225 页）。
它在 partinfo 里的坐标（\chref{ch:info-climate-change-adaptation}）：方法线 panel line 的第三站 climate penalty approach。它回答第二站 nonlinear panel approach 留下的问题：proposition 只说二次面板的系数里含有 long-run 成分，没有说含多少。它给出的增广回归与 cross-sectional line 上的 long differences 做的是同一件事，即比较 short-run 与 long-run response，但不需要各 location 之间长期趋势的差异。与 \chref{ch:hsiang-2016} 的关系：Hsiang 的 Envelope 论证给出两种效应在边际处相等的条件，论文的结果说明离开这一点之后面板系数偏离多少（第~1217 页）。与两篇经验研究的关系：\chref{ch:deschenes-greenstone-2007} 自认面板估计不含全部 long-run 调整，\chref{ch:burke-hsiang-miguel-2015} 的二次面板系数如何解读没有定论，论文的偏差分解对两者都给出了可计算的回答。它的第一篇应用是本 part 的 Application 1（\chref{ch:da-et-al-2024}）。
% 坐标表述：2026-11-19 按主笔意见随 partinfo 两条线的方法线重写（原批注「暂不修改」已解除）。
骨架表（所据版本：AJAE 103(4) 已刊版；页码为期刊页码）：

\begin{table}[htbp]
\centering\small
\begin{tabular}{@{}p{5.4cm}p{1.6cm}p{6.0cm}@{}}
\toprule
原文章节 & 页码 & 讲什么 \\
\midrule
Abstract、Introduction & 1207--1210 & 两条方法与主张；研究设计的定位 \\
A Simple Model of Long-Run Climatic Adaptation & 1210--1214 & climate penalty DGP 与行为推导（本讲义第二节精讲） \\
Asymptotic Bias with a Linear Long-Run Effect & 1215 & 式 (3)--(5)：偏度机制 \\
Asymptotic Bias with a Nonlinear Long-Run Effect & 1215--1220 & 一般偏差（式 (7)--(9)）、无 skewness 情形（式 (10)--(13)）、一般分解（式 (14)(15)） \\
Consistent Estimation of Short-Run and Long-Run Responses and Counterfactual Analysis & 1220--1222 & 增广回归（式 (16)）与 counterfactual（反事实，式 (17)） \\
Trends in Historical Climate & 1222--1224 & climate 有趋势时结论如何变化；Table~1 \\
Empirical Applications & 1224--1235 & Bias Decomposition（Table~2）；Joint Estimation（Table~3--4、Figures~5--8）；Vulnerability of Panel Estimates to Omitted Variables（式 (18)(19)） \\
Conclusion & 1235--1236 & 结论与政策含义 \\
\bottomrule
\end{tabular}
\end{table}

\section{预备知识}

三件计量预备知识，用前备齐：\emph{plim} 记号（估计量按概率收敛的极限）；\index{convex combination}convex combination（凸组合，权重非负且和为一的加权平均）；\index{within estimator}within estimator（组内估计量，用组内去均值 variation 做 OLS 的 fixed effects 估计量）。论文的渐近框架是固定 location 数 $I$、时期数 $T\to\infty$：``Our asymptotic results are derived for a fixed set of locations and a large number of time periods''\citep[第~1209 页]{merel2021climate}；climate 带趋势的一般情形放在在线附录 F，正文取 stationary 情形——其经验理由是 climate 趋势在论文所用的数据里占比很小（第~1223--1224 页，见「结果与解释」节）。
\section{A Simple Model of Long-Run Climatic Adaptation}
\label{sec:cca-model}
本节基于论文第二节与 McIntosh 和 Schlenker（2006），讲清 climate penalty panel model 的 identification 逻辑：要识别什么、依赖哪条 assumption、估得的系数怎么读，最后给出几何图景。论文自述这是一个 reduced-form 的 long-run adaptation 模型（第~1210 页）：它先写出主体的行动选择问题，再由最优行动推出只含 weather 与 climate 的 reduced-form DGP，行动本身不出现在估计方程里。

\subsection{Identification model：数据、方程与 Assumption 1}

数据是 $I$ 个 location（如县）乘 $T$ 个时期（如年）的面板。每个观测包含三个要素：产出 $y_{it}$（论文原例为对数农场利润，经验部分用作物单产）；\index{weather}weather $x_{it}$，即 location $i$ 在时期 $t$ 的 weather 实现值（growing season 的平均温度、降水等），逐年波动；以及 location 的 \index{climate}climate $\mu_i$，实践中取 weather 的长期均值。weather 与 climate 的区分沿用本 part 的约定（\chref{ch:info-climate-change-adaptation}）。

要估计的方程是
\begin{equation}
y_{it} = \alpha_i + \beta_1 x_{it} + \beta_2 x_{it}^2 + \beta_3 (x_{it}-\mu_i)^2 + \eps_{it},
\qquad \beta_3 \le 0 .
\label{eq:cca-dgp}
\end{equation}
这是论文的式 (1)（第~1211 页）：前两项是 weather 的二次响应，二次型允许 weather 与产出的关系非单调（$\beta_2<0$ 时呈穹形，第~1211 页）；第三项度量「当年 weather 偏离该 location 自身 climate」带来的损失，论文把它称为 \index{climate penalty}climate penalty term（第~1213 页）。本讲义把式~\ref{eq:cca-dgp} 称作 climate penalty model，\chref{ch:da-et-al-2024} 所读的应用论文也用这个名字。DGP（data-generating process，数据生成过程）指的是这组方程所描述的「产出如何由 weather、climate 与 adaptation 行为生成」的机制，而不是某一种数据：它既决定短期观测 $y_{it}$ 的形态，也蕴含长期 counterfactual。该设定与 McIntosh 和 Schlenker（2006）提出的模型相同（第~1211 页；\citealp{mcintosh2006identifying}）。

\textbf{Assumption 1（论文未编号的展示式，第~1210 页；论文此后称之为 the stationarity assumption）。}观测期内每个 location 的 climate 是 stationary 的：weather 是以该 location 的 climate 为中心的随机变量，
\begin{equation}
\E[x_{it} \mid i] = \mu_i ,
\label{eq:cca-climate}
\end{equation}
即给定 location，其 climate $\mu_i$ 在观测期内不变，weather 只是围绕它波动。这条 assumption 把 adaptation 的位置说清楚了：adaptation 指给定 location 的主体对「weather 偏离本地 climate」的既有适应，penalty term $(x_{it}-\mu_i)^2$ 正是这一偏离的函数。它与另一种读法要有分工：若把 adaptation 理解为「climate 变化之后主体的再调整」，那是 long-run adaptation，在数据里只体现于跨 location 的 climatic differences（见下文 counterfactual 与几何部分）；Assumption 1 要求观测期内没有这种 climate 漂移，观测到的只是对既有 climate 的适应结果。由此也能回答一个常见疑问：long-run adaptation 并不是「藏在 location fixed effects $\alpha_i$ 里」的——$\alpha_i$ 只是截距，吸收不随时间变化的产出水平差异；long-run 信息在 $\mu_i$ 里，并通过 penalty term 进入方程，而把 $\mu_i$ 解释成 climate 靠的正是 stationarity。

观测期就是样本实际覆盖的年份，论文没有另设一个「climate 不变年限」，Assumption 1 是对这段观测期整体的要求。样本期越长，climate 趋势越可能显形，论文因此把带趋势 climate 的推广放到在线附录 F，并用数据说明所用样本中 climate 趋势占比很小（第~1223--1224 页，见「结果与解释」节）。估计层面另需通常的强外生性（strong exogeneity），论文只作文字陈述（第~1214 页），写成公式是 $\E[\eps_{it}\mid \alpha_i,\mu_i,x_{i1},\dots,x_{iT}]=0$（本讲义的形式化）。时间效应方面：只要 climate 是 stationary 的，weather 方差 $\E[(x_{it}-\mu_i)^2]$ 不随时间改变，被漏掉的 penalty term 与任何时间 fixed effects 渐近不相关（第~1214 页）。

\begin{intuition}[optimizing agents：penalty term 的行为来源]
penalty term 不是额外添加的回归项，论文把它从主体的优化问题中推了出来（第~1211 页）。设 location $i$ 的主体选择行动集 $\boldsymbol{\xi}_i\in\R^L$（品种、种植日期、灌溉等），产出由函数 $\phi_i$ 给出：
\begin{equation*}
y_{it}=\phi_i(x_{it},\boldsymbol{\xi}_i,\eps_{it})
=\alpha_i+\beta_1 x_{it}+\beta_2 x_{it}^2+\beta_3\bigl(x_{it}-\omega_i(\boldsymbol{\xi}_i)\bigr)^2+\eps_{it},
\end{equation*}
其中 $\omega_i:\R^L\to X_i$ 为满射；条件于 weather $x_{it}$，产出当且仅当 $\omega_i(\boldsymbol{\xi}_i)=x_{it}$ 时最大。主体选 $\boldsymbol{\xi}_i$ 以最大化期望产出：
\begin{align*}
\boldsymbol{\xi}_i&\in\operatorname{argmax}\,\E[\phi_i(x_{it},\boldsymbol{\xi}_i,\eps_{it})]
=\operatorname{argmin}\,\E\bigl[(x_{it}-\omega_i(\boldsymbol{\xi}_i))^2\bigr]\\
&=\operatorname{argmin}\,\E\bigl[(x_{it}-\mu_i+\mu_i-\omega_i(\boldsymbol{\xi}_i))^2\bigr]
=\operatorname{argmin}\,(\mu_i-\omega_i(\boldsymbol{\xi}_i))^2
=\omega_i^{-1}\bigl(\{\mu_i\}\bigr).
\end{align*}
倒数第二步把平方展开：$\E[(x_{it}-\mu_i)^2]$ 与行动无关，交叉项 $2\E[x_{it}-\mu_i](\mu_i-\omega_i(\boldsymbol{\xi}_i))$ 为零，用的正是 Assumption 1。于是最优行动使 $\omega_i(\boldsymbol{\xi}_i)=\mu_i$，代回即得式~\ref{eq:cca-dgp}（论文式 (1)，第~1211 页）。行动下标只有 $i$ 没有 $t$：行动不针对逐年 weather 调整，论文称之为 committed actions（第~1211 页）。

短期与长期的差别不在「是否理性」，而在哪些行动来得及调整。论文举了两个例子（第~1212--1213 页）：其一引自 \citet{auffhammer2014empirical}，一年的干旱不值得为此修建灌渠，但在干旱的 climate 里修灌渠可能是值得的；其二是论文自己的类比，异常热浪不值得为此安装整套空调系统，但会促使人买一台吊扇。长期投入只对 climate 响应，在同一函数形式里由 penalty term 体现。

论文还指出式~\ref{eq:cca-dgp} 容许一个更结构化的读法（第~1211--1212 页）：在 Schlenker（2017）的品种选择框架里，$\beta_3$ 可以依赖内生选择的品种指数 $\gamma_i$，农民依据 climate 的均值与方差选品种，penalty term 中 $\gamma_i$ 取代 $\mu_i$，此时 penalty 的最低点未必落在 $\mu_i$；Schlenker 发现这层灵活性对美国玉米单产的 long-run 系数估计没有实质改变 \citep{schlenker2017interannual}。
\end{intuition}

\subsection{Interpretation：short-run 拟合、long-run counterfactual 与 adaptation 系数}

式~\ref{eq:cca-dgp} 拟合的是 \index{short-run response}short-run response：在 $\alpha_i$ 与 $\mu_i$ 给定、行动已按本地 climate 安排的前提下，产出随 weather 实现的响应（第~1211 页），
\begin{equation}
\E[y_{i}^{\mathrm{SR}} \mid \alpha_i, \mu_i, x_i] = \alpha_i + \beta_1 x_i + \beta_2 x_i^2 + \beta_3 (x_i-\mu_i)^2 .
\label{eq:cca-sr}
\end{equation}
但评估 climate change 关心的对象是 \index{long-run response}long-run response：当 climate 本身变化、主体完成全部再适应之后，penalty 消失（长期里所有行动都可随环境条件调整，第~1211 页），
\begin{equation}
\E[y_{i}^{\mathrm{LR}} \mid \alpha_i, x_i] = \alpha_i + \beta_1 x_i + \beta_2 x_i^2 .
\label{eq:cca-lr}
\end{equation}
论文把基准对象定为 long-run 系数 $\boldsymbol{\beta}^{\mathrm{LR}}=(\beta_1,\beta_2)$（第~1214 页）。climate 从 $\mu_0$ 变到 $\mu_1$ 的长期 counterfactual 损害是 $\Delta^{\mathrm{LR}}=\beta_1(\mu_1-\mu_0)+\beta_2(\mu_1^2-\mu_0^2)$（第~1213 页）；这是在 weather 均值处评价的产出变化，不是期望产出的变化，后者还含 weather 方差的项，见「一致估计」节（论文脚注~13，第~1213 页）。

关键限制：long-run response 在数据里观测不到。每个观测只是某条 short-run response 曲线上的一个点，$\mu_i$ 不变，主体也不会为单年 weather 重做长期投入（committed actions）。因此 long-run counterfactual 不能单独估计，只能经由式~\ref{eq:cca-dgp} 的短期识别反推：短期方程里的 $\beta_1,\beta_2$ 就是 long-run response 的系数。

由此 adaptation 由 $\beta_3$ 承载。$\beta_3\le 0$ 度量「weather 偏离本地 climate」的 penalty 强度；$\beta_3=0$ 即没有 climatic adaptation 的情形（第~1211 页），此时 short-run 与 long-run response 重合；$|\beta_3|$ 越大，同一偏离造成的损失越大，short-run 曲线相对 long-run 基准下沉越深。注意 $\beta_3$ 度量的是 penalty 的强度，而不是 adaptation 抵消了多少损害——后者要由 counterfactual 中 short-run 与 long-run 效应之差给出（见「一致估计」节式~\ref{eq:cca-cf}）。

先算一个例子给读者看。设甲县 climate 偏冷 $\mu=16\,^\circ$C，乙县 climate 偏热 $\mu=24\,^\circ$C，penalty 系数取 $\beta_3=-0.01$（对数产出的示意参数，为本讲义自设，不是原文数字）。growing season 平均温度实现 $x=21\,^\circ$C 时：甲县 penalty $-0.01\times(21-16)^2=-0.25$ 个对数点，乙县 penalty $-0.01\times(21-24)^2=-0.09$ 个对数点。同一场 weather，climate 更接近它的乙县损失更小。把 weather 换成 $x=19\,^\circ$C 再算一遍：甲县 $-0.01\times(19-16)^2=-0.09$，乙县 $-0.01\times(19-24)^2=-0.25$，损失恰好对调。这个两县例子在后面讲偏差时还要用。

\begin{pitfall}[把 climate 均值当成暴露度的完备刻画]
\label{pit:cca-mu}
式~\ref{eq:cca-dgp} 有两个限制，转述与解释系数时不可丢掉。其一是论文自己写明的，penalty 只看绝对距离、不看方向（第~1213 页）：上面两县例子里，$x=21\,^\circ$C 对乙县的 3\,$^\circ$C 偏差与 $x=19\,^\circ$C 对甲县的 3\,$^\circ$C 偏差被同等对待；论文自带的例子是冷热两个 location 都暴露在两地 climate 均值的 weather 下，penalty 相同。若冷偏离与热偏离的后果不对称，二次 penalty 就把两种后果压平了。其二，均值未必是 adaptation 所对准的基准。已刊版的相关说法有两处：climate 在这里指 weather 的均值，而广义的 climate 是 weather 的分布，可以由一阶矩以外的量刻画（第~1210 页）；在 Schlenker（2017）的框架下，penalty 的最低点未必落在 climate 均值处（第~1211--1212 页）。本讲义原稿所举的「climate 分布双峰时均值刻画能力很弱」一例出自工作稿第一版（第 5 页），已刊版没有保留，本仓亦无该工作稿，待复核，引用时不可归于已刊版。误用的形态是把「对均值的偏离」当成「对实际暴露的偏离」来解释 penalty term 系数。
\end{pitfall}

把短期与长期两个对象接起来有两条直接的警告。第一条是论文在第~1213 页给出的：把式~\ref{eq:cca-dgp} 展开，可以改写成含 climate 与 weather 交互项 $\mu_i x_{it}$ 的方程（论文式 (2)），其中 $x_{it}^2$ 的系数是 $\beta_2+\beta_3$；若把这一改写式当作 long-run 关系直接计算 counterfactual，得到的是 $\beta_1(\mu_1-\mu_0)+(\beta_2-\beta_3)(\mu_1^2-\mu_0^2)$，只要 $\beta_3\neq0$ 就与真实的 $\Delta^{\mathrm{LR}}$ 不同，原因是改写式的截距里含 $\beta_3\mu_i^2$，它随 climate 变化（第~1213--1214 页）。第二条关于 naive 路线，即漏掉 penalty term 的二次面板：它估得的曲率是 $\beta_2+\bar\theta\beta_3$ 而非 $\beta_2$，$\bar\theta$ 是下一节推导的权重（第~1217 页）。「panel approach 估计到的到底是什么、偏向 long-run 多少」，是下文「渐近偏差分解」与「一致估计」两节的任务。

\begin{ident}[带 climate penalty 的二次面板：可整体搬用的设定]
\begin{align*}
\text{DGP：}\quad & y_{it} = \alpha_i + \beta_1 x_{it} + \beta_2 x_{it}^2 + \beta_3 (x_{it}-\mu_i)^2 + \eps_{it}, \qquad \beta_3 \le 0 \\
\text{long-run response：}\quad & \E[y_{i}^{\mathrm{LR}} \mid \alpha_i, x_i] = \alpha_i + \beta_1 x_i + \beta_2 x_i^2 \\
\text{short-run response：}\quad & \E[y_{i}^{\mathrm{SR}} \mid \alpha_i, \mu_i, x_i] = \alpha_i + \beta_1 x_i + \beta_2 x_i^2 + \beta_3 (x_i-\mu_i)^2 \\
\text{identification：}\quad & \E[x_{it}\mid i]=\mu_i,\qquad \E[\eps_{it} \mid \alpha_i, \mu_i, x_{i1}, \dots, x_{iT}] = 0
\end{align*}
配套的口径要求：climate $\mu_i$ 取 weather 的长期均值且样本期内不变；penalty term 以「weather 对本 location climate 的偏差」为自变量，与 weather 本身的三项（截距、一次、二次）并存。强外生性的公式为本讲义的形式化（论文只作文字陈述，第~1214 页）。
\end{ident}

\subsection{几何理解：long-run 包络线与 short-run 曲线族}

short-run 与 long-run 的几何形态见图~\ref{fig:cca-lrsr}，依据已刊版图 1（第~1212 页）重绘。取 $\beta_2<0$ 的情形，long-run response 呈穹形（第~1211 页）；$\beta_3\le0$ 使每个 location 的 short-run response 曲线都不高于 long-run response，并在自己的 climate $\mu_i$ 处与之相切，因为 $x=\mu_i$ 时 penalty 恰为零、两条曲线的斜率也在此相等（第~1211 页）。把各 location 的 $\alpha_i$ 画成相同之后，long-run response 恰是 short-run response 曲线族的上包络（第~1211 页）。原图还在各 location 的 short-run response 曲线附近画了样本数据点与箭头，重绘从略。weather 离本地 climate 越远，虚线沉得越深：short-run response 对「不熟悉的 weather」更敏感。「naive 面板模型拟合出的曲线落在 long-run response 与 short-run response 之间的哪里」，就是这条几何图景提出、论文后续各节回答的问题。

\begin{figure}[htbp]
  \centering
  \begin{tikzpicture}
    \begin{axis}[
        width=0.82\textwidth, height=6.4cm,
        xlabel={weather $x$}, ylabel={期望产出 $y$},
        axis lines=left,
        xmin=0, xmax=4.4, ymin=-2.4, ymax=5.2,
        xtick={1,2,3}, xticklabels={$\mu_1$,$\mu_2$,$\mu_3$},
        ytick=\empty,
        tick label style={font=\small},
        label style={font=\small},
        legend pos=north east,
        legend style={font=\small, draw=none, fill=none},
        samples=120]
      \addplot[cMain, very thick, domain=0:4] {4*x - x^2};
      \addlegendentry{long-run response（式~\ref{eq:cca-lr}）}
      \addplot[cBranch, thick, dashed, domain=0:3.15] {4*x - x^2 - (x-1)^2};
      \addlegendentry{short-run response（式~\ref{eq:cca-sr}）}
      \addplot[cBranch, thick, dashed, domain=0.25:3.75] {4*x - x^2 - (x-2)^2};
      \addplot[cBranch, thick, dashed, domain=1.05:4] {4*x - x^2 - (x-3)^2};
      \addplot[only marks, mark=*, mark size=1.7pt, black]
        coordinates {(1,3) (2,4) (3,3)};
      \node[lbl, anchor=south] at (axis cs:2,4.15) {切点：$x=\mu_i$ 处 penalty 为零};
    \end{axis}
  \end{tikzpicture}
  \caption{long-run response 与三个 location 的 short-run response（依据 Mérel 与 Gammans（2021）已刊版图~1 重绘，原图在第~1212 页；原图另含样本数据点与箭头，此处从略）。参数为本讲义自选（$\beta_1=4$、$\beta_2=-1$、$\beta_3=-1$、各 location $\alpha_i$ 相同），只示意曲线形状与切点位置。$\beta_2<0$ 时 long-run response 呈穹形；$\beta_3<0$ 时它是 short-run response 曲线族的上包络，每个 location 的 short-run response 在其 climate $\mu_i$ 处与 long-run response 相切，weather 离本地 climate 越远，short-run 产出相对 long-run 基准的 penalty 越重。若各 location $\alpha_i$ 不同，各曲线竖直平移，切点的横坐标不变（第~1211 页）。}
  \label{fig:cca-lrsr}
\end{figure}

本节简化模型的结论自带适用范围：式~\ref{eq:cca-dgp} 是论文为可解析性选定的特定函数形式，「偏差由 skewness 驱动」「convex combination」这些结论都是在它之上推出来的。论文对二次设定适用面的说明有两处：以 weather 的非线性变换（如 degree-days）作回归元的设定可以改写进线性或二次框架（第~1208 页脚注~3）；经验部分所重访的研究并非都用二次设定，权重参数的计算是对「估计离 long-run response 有多近」的一个大致判断（第~1224 页）。换成别的非线性设定是否仍有同型的闭式结论，论文没有给出。两县例子（$\mu=16\,^\circ$C 与 $\mu=24\,^\circ$C）后续小节继续沿用。

\section{渐近偏差分解：naive 面板系数的构成}
\label{sec:cca-bias}

% 主笔批注（2026-11-06）：section 导读没有问题；subsection 的推导逻辑改为
% 按 omitted variable bias 的推导过程计算偏差公式与 convex combination 两式
% （当时编号 5.7、5.9，即 eq:cca-biaslinear 与 eq:cca-combo），正文写清推导逻辑，
% 用卡片写具体计算过程。2026-11-19 两个小节均已照此改写。
本节完成「A Simple Model of Long-Run Climatic Adaptation」一节留下的任务：在式~\ref{eq:cca-dgp} 之下，计算漏掉 penalty term 的 naive 二次面板的估计量收敛到什么。先看论文自己的一句话总结：``Said differently, panel estimates from the quadratic weather-outcome relationship can be thought of as a weighted average of short- and long-run responses, with the weight on the long-run coefficients increasing with the share of overall weather variation attributable to cross-sectional (i.e., climatic) differences''\citep[第~1209 页]{merel2021climate}。本节其余内容就是把这句话拆开：权重是什么、由什么决定、偏差朝哪边。

两个小节用同一条推导逻辑。naive 回归把 $w_{it}=(x_{it}-\mu_i)^2$ 漏在误差项里，所以它的 within estimator 的概率极限等于真实系数加上一项 omitted variables bias：$\beta_3$ 乘以「$w_{it}$ 对已含回归元做组内回归的系数」。剩下的工作只是把后者用 weather 的矩算出来。记号随论文（第~1215 页）：$\ddot{x}_{it}$ 是 $x_{it}$ 减去本 location 的时间均值，$T\to\infty$ 时即 $x_{it}-\mu_i$；location $i$ 的二阶、三阶矩记作 $M^i_{\ddot{x}^2}$ 与 $M^i_{\ddot{x}^3}$，本讲义为书写方便另记 $\sigma_i^2\equiv M^i_{\ddot{x}^2}$、$\mu_{3i}\equiv M^i_{\ddot{x}^3}$。

\subsection{线性 long-run effect：偏度机制（式 (3)--(5)）}

先取 $\beta_2=0$ 的特例。DGP 为 $y_{it} = \alpha_i + \beta_1 x_{it} + \beta_3 (x_{it}-\mu_i)^2 + \eps_{it}$（式 (3)），naive 模型为 $y_{it} = a_i + b_1 x_{it} + e_{it}$（式 (4)），用 within estimator 估计 $b_1$，$T\to\infty$。这里只有一个已含回归元，omitted variables bias 公式是一个比值：分子是 $\ddot{x}_{it}$ 与被漏变量的协方差，分母是 $\ddot{x}_{it}$ 的方差，各 location 合并。算出来是
\begin{equation}
\mathrm{plim}\,\widehat{b}_1 - \beta_1 \;=\; \beta_3\,\frac{\sum_i \mu_{3i}}{\sum_i \sigma_i^2},
\qquad
\mu_{3i} \equiv \E\bigl[(x_{it}-\mu_i)^3\bigr],\quad
\sigma_i^2 \equiv \E\bigl[(x_{it}-\mu_i)^2\bigr],
\label{eq:cca-biaslinear}
\end{equation}
即偏差等于「penalty 系数 $\beta_3$」乘「各 location weather 三阶中心矩之和除以 weather 方差之和」；$\mu_{3i}$ 未除以 $\sigma_i^3$，是未归一化的 \index{skewness}skewness（偏度）。这就是论文的式 (5)（第~1215 页，以矩记号 $M^i_{\ddot{x}^3}$、$M^i_{\ddot{x}^2}$ 写出），分子为什么是三阶矩由其脚注点明：``The covariance between the included regressor $x_{it}$ and the omitted variable $\ddot{x}_{it}^2$ is equal to $\E[\ddot{x}_{it}^3]$''\citep[第~1215 页脚注~18]{merel2021climate}。化简过程见本小节末的推导卡片。

式~\ref{eq:cca-biaslinear} 连带三条读法。其一，偏差大小与被忽略的 adaptation 成比例：``the size of the bias is proportional to the extent of adaptation being ignored in the estimated model ($\beta_3$)''\citep[第~1215 页]{merel2021climate}。其二，方向完全由 skewness 决定：``the bias is entirely driven by the skewness in the weather data conditional on location''\citep[第~1215 页]{merel2021climate}；weather 右偏（$\mu_{3i}>0$）时 $\beta_3<0$ 把斜率压得更负，增温的损害被高估，weather 左偏则反过来（第~1215 页）。其三，对称分布无偏差：对称时 $\mu_{3i}=0$，``larger values of $\ddot{x}_{it}$ are not systematically associated with larger values of $\ddot{x}_{it}^2$''\citep[第~1215 页]{merel2021climate}——被漏变量与回归元的协方差为零。

用两县例子算一遍（示意参数沿用第二节：$\beta_3=-0.01$；本讲义推算，读者可验算）。设甲县 weather 相对 $\mu=16\,^\circ$C 的偏离只有两个取值：$-1\,^\circ$C（概率 0.8）与 $+4\,^\circ$C（概率 0.2），右偏；乙县偏离对称取 $\pm2\,^\circ$C（各 0.5）。甲县 $\sigma^2=0.8\times1+0.2\times16=4$、$\mu_3=0.8\times(-1)+0.2\times64=12$；乙县 $\sigma^2=4$、$\mu_3=0$。合计 $\sum\mu_3=12$、$\sum\sigma^2=8$，式~\ref{eq:cca-biaslinear} 给 $\mathrm{plim}\,\widehat{b}_1-\beta_1=-0.01\times1.5=-0.015$：右偏县的热浪年份里 penalty 与大的正偏离同时出现，naive 斜率把这部分 penalty 误记成线性温度效应，增温损害被多算。论文的 Figure~2（第~1216 页）画的是同一机制的连续版：右偏分布下，估计出的直线比线性的 long-run response 更陡。

\begin{intuition}[线性情形：omitted variables bias 公式逐步化简]
组内去均值后，naive 回归把 $w_{it}=(x_{it}-\mu_i)^2=\ddot{x}_{it}^2$ 漏在误差项里，omitted variables bias 公式给出
\begin{equation}
\mathrm{plim}\,\widehat{b}_1-\beta_1=\beta_3\,\frac{\sum_i\Cov_i(\ddot{x}_{it},w_{it})}{\sum_i\Var_i(\ddot{x}_{it})},
\label{eq:cca-ovb}
\end{equation}
下标 $i$ 表示在 location $i$ 内对时间取矩。分子对单个 location 化简：
\begin{align}
\Cov_i\bigl(\ddot{x}_{it},\ddot{x}_{it}^{2}\bigr)
&=\E\bigl[\ddot{x}_{it}\cdot\ddot{x}_{it}^{2}\bigr]-\E[\ddot{x}_{it}]\,\E\bigl[\ddot{x}_{it}^{2}\bigr] \notag\\
&=\E\bigl[\ddot{x}_{it}^{3}\bigr]-0\cdot\sigma_i^2
=\mu_{3i},
\label{eq:cca-moment}
\end{align}
第一步是协方差定义，第二步用 Assumption 1（$\E[\ddot{x}_{it}]=0$）；分母即 $\sigma_i^2$。对各 location 求和，$\mathrm{plim}\,\widehat{b}_1-\beta_1=\beta_3\sum_i\mu_{3i}\big/\sum_i\sigma_i^2$，即式~\ref{eq:cca-biaslinear}（论文式 (5)，第~1215 页）。
\end{intuition}

\subsection{非线性 long-run effect：无偏度时的凸组合（式 (7)--(12)）}

$\beta_2\neq 0$ 时 naive 模型（论文式 (6)：$y_{it}=a_i+b_1x_{it}+b_2x_{it}^2+e_{it}$）与 DGP（式~\ref{eq:cca-dgp}）都是二次的。推导逻辑不变，只是已含回归元从一个变成两个（$x_{it}$ 与 $x_{it}^2$），omitted variables bias 公式从一个比值变成一个 $2\times2$ 的矩阵式：
\begin{equation}
\mathrm{plim}\,\widehat{\mathbf{b}}-\boldsymbol{\beta}^{\mathrm{LR}}
=\beta_3\Bigl[\sum_i \mathbf{Q}_i\Bigr]^{-1}\sum_i \mathbf{c}_i ,
\label{eq:cca-ovb2}
\end{equation}
其中 $\mathbf{Q}_i$ 是 $(x_{it},x_{it}^2)$ 在 location $i$ 内的协方差矩阵，$\mathbf{c}_i$ 是 $(x_{it},x_{it}^2)$ 与被漏变量 $w_{it}$ 在 location $i$ 内的协方差向量（记号 $\mathbf{Q}_i$、$\mathbf{c}_i$ 为本讲义所设）。把矩代进去要用到一个新的量：偏离平方 $\ddot{x}_{it}^2$ 在 location 内的方差，本讲义记作 $v_i\equiv\Var_i(\ddot{x}_{it}^2)$。论文对应的记号是 $M^i_{(\ddot{x}^2-\ddot{x}^2)^2}\equiv\mathrm{plim}\,T^{-2}\sum_{s,t}(\ddot{x}_{is}^2-\ddot{x}_{it}^2)^2$（第~1216 页）；对时期的无序对求和时它等于 $v_i$，本讲义的推导结果与论文式 (12) 正是在这一读法下逐项吻合。论文指出，weather 只取两个等概率的值时这个矩为零，否则为正（第~1216 页）。

一般情形下，marginal effect 的偏差依赖评价点 climate $\mu$。论文定义 $\mathrm{Bias}(\mu)=\mathrm{plim}\,\widehat{b}_1-\beta_1+2\mu(\mathrm{plim}\,\widehat{b}_2-\beta_2)$，结果是
\begin{equation}
\mathrm{Bias}(\mu) \;=\; -\,\frac{2\beta_3\,N}{D} ,
\label{eq:cca-biasmu}
\end{equation}
其中 $N$、$D$ 由各 location 的 $M^i_{\ddot{x}^2}$、$M^i_{\ddot{x}^3}$、$M^i_{(\ddot{x}^2-\ddot{x}^2)^2}$、climate 水平 $\mu_i$ 与 climatic differences $(\mu_i-\mu_j)^2$ 构成（式 (7)--(9)，第~1216 页）。与线性情形相同，含三阶矩的项使偏差的符号一般不能确定；与线性情形不同的是，偏差不再完全由 skewness 驱动（第~1216--1217 页）。所以论文转入 ``Bias When the Weather Distribution Is Not Skewed''\citep[第~1217 页]{merel2021climate}（$M^i_{\ddot{x}^3}=0\ \forall i$）的专项分析：此时式~\ref{eq:cca-biasmu} 简化为式 (10)，within estimator 收敛到一个 convex combination（式 (11)）
\begin{equation}
\mathrm{plim}\,\widehat{\mathbf{b}} \;=\; (1-\bar\theta)\,\boldsymbol{\beta}^{\mathrm{LR}} \;+\; \bar\theta\sum_i \lambda_i\,\boldsymbol{\beta}^{\mathrm{SR}}_{i},
\qquad
\boldsymbol{\beta}^{\mathrm{SR}}_{i} = \begin{pmatrix} \beta_1 - 2\beta_3\mu_i \\ \beta_2+\beta_3 \end{pmatrix},
\label{eq:cca-combo}
\end{equation}
其中 $\boldsymbol{\beta}^{\mathrm{SR}}_{i}$ 是把式~\ref{eq:cca-sr} 展开后 location $i$ 的 short-run 一次项与二次项系数；\index{weight parameter}weight parameter（权重参数）$\lambda_i = \sigma_i^2\big/\sum_j \sigma_j^2$ 可以解读为``the share of overall time-series weather variation attributable to location $i$''\citep[第~1217 页]{merel2021climate}；$\bar\theta$ 由式 (12) 给出，用本讲义的记号写是
\begin{equation}
\bar\theta=\frac{\sum_i\sigma_i^2\,\sum_i v_i}{\sum_i\sigma_i^2\,\sum_i v_i+4\sum_{i<j}(\mu_i-\mu_j)^2\sigma_i^2\sigma_j^2}\;\in[0,1] .
\label{eq:cca-theta}
\end{equation}
$\bar\theta$ 是 short-run 系数上的权重，$1-\bar\theta$ 是 long-run 系数上的权重。分子分母共有的一项只含各 location 内部的时间序列 variation；分母多出的一项是成对的 climatic differences。横截面差异相对组内 variation 越大，$\bar\theta$ 越接近零，naive 系数越接近 long-run 系数。由式~\ref{eq:cca-combo} 的第二个分量直接读出 $\mathrm{plim}\,\widehat{b}_2-\beta_2=\beta_3\bar\theta$（第~1217 页）。化简过程见下面的卡片。

几何解读（论文 Figure~3，第~1218 页）：记组合 climate $\bar{\mu} = \sum_i \lambda_i \mu_i$（第~1217 页），naive 估计的曲线恰在 $\bar\mu$ 处与 long-run response 斜率相同——评价点取在 $\bar{\mu}$ 则 marginal effect 无偏；评价点 $\mu<\bar{\mu}$ 时偏差为正，$\mu>\bar{\mu}$ 时偏差为负（第~1218 页）。这直接给出一个加总警告：各 location 的偏差有正有负，加总后的净偏差取决于 $\lambda_i$ 的结构、各地 climate 变化的幅度与加总权重，可以落在任一方向（第~1218 页）。有 skewness 的一般情形下，分解的形式不变，权重与切点换成 $\tilde\theta$（式 (15)）与 $\tilde\mu$（式 (14)），二者都含三阶矩（第~1219--1220 页）。

\begin{intuition}[二次情形：由 omitted variables bias 公式到 convex combination]
记 $d\equiv\ddot{x}_{it}=x_{it}-\mu_i$，在 location $i$ 内 $\E[d]=0$、$\E[d^2]=\sigma_i^2$、$\E[d^3]=\mu_{3i}$、$\Var(d^2)=v_i$。把 $x_{it}^2$ 写成 $\mu_i^2+2\mu_i d+d^2$，逐项计算。

第一步，已含回归元的组内协方差矩阵。$\Var_i(x)=\sigma_i^2$；$\Cov_i(x,x^2)=2\mu_i\sigma_i^2+\mu_{3i}$；$\Var_i(x^2)=4\mu_i^2\sigma_i^2+4\mu_i\mu_{3i}+v_i$。

第二步，与被漏变量 $w=d^2$ 的协方差。$\Cov_i(x,w)=\mu_{3i}$（即式~\ref{eq:cca-moment}）；$\Cov_i(x^2,w)=2\mu_i\mu_{3i}+v_i$。

第三步，令 skewness 为零（$\mu_{3i}=0$）并对 location 求和。记 $S=\sum_i\sigma_i^2$、$A=\sum_i\mu_i\sigma_i^2$、$B=\sum_i\mu_i^2\sigma_i^2$、$V=\sum_i v_i$，则
\begin{equation*}
\sum_i\mathbf{Q}_i=\begin{pmatrix} S & 2A\\ 2A & 4B+V\end{pmatrix},\qquad
\sum_i\mathbf{c}_i=\begin{pmatrix}0\\ V\end{pmatrix}.
\end{equation*}

第四步，代入式~\ref{eq:cca-ovb2}。行列式为 $S(4B+V)-4A^2=SV+4(SB-A^2)$，于是
\begin{align*}
\mathrm{plim}\,\widehat{\mathbf{b}}-\boldsymbol{\beta}^{\mathrm{LR}}
&=\frac{\beta_3}{SV+4(SB-A^2)}\begin{pmatrix}4B+V & -2A\\ -2A & S\end{pmatrix}\begin{pmatrix}0\\ V\end{pmatrix}\\
&=\beta_3\,\frac{SV}{SV+4(SB-A^2)}\begin{pmatrix}-2A/S\\ 1\end{pmatrix}.
\end{align*}

第五步，认出各项。$SB-A^2=\sum_{i<j}(\mu_i-\mu_j)^2\sigma_i^2\sigma_j^2$（把两个和式展开后两两配对即得），所以前面的分式就是式~\ref{eq:cca-theta} 的 $\bar\theta$；$A/S=\sum_i\lambda_i\mu_i=\bar\mu$。结果写成
\begin{equation*}
\mathrm{plim}\,\widehat{b}_1=\beta_1-2\beta_3\bar\theta\bar\mu,\qquad
\mathrm{plim}\,\widehat{b}_2=\beta_2+\bar\theta\beta_3 .
\end{equation*}
而 $(1-\bar\theta)\boldsymbol{\beta}^{\mathrm{LR}}+\bar\theta\sum_i\lambda_i\boldsymbol{\beta}^{\mathrm{SR}}_i$ 的两个分量正是 $\beta_1-2\beta_3\bar\theta\sum_i\lambda_i\mu_i$ 与 $\beta_2+\bar\theta\beta_3$，即式~\ref{eq:cca-combo}（论文式 (11)(12)，第~1217 页）。marginal effect 的偏差随之为 $\mathrm{Bias}(\mu)=2\beta_3\bar\theta(\mu-\bar\mu)$，与论文式 (10) 一致：$\beta_3<0$ 时，$\mu<\bar\mu$ 处为正，$\mu>\bar\mu$ 处为负。以上推导为本讲义按 omitted variables bias 公式所作，记号 $S$、$A$、$B$、$V$ 为本讲义所设；论文正文只给出结果。
\end{intuition}

用两县例子把权重算出来（本讲义推算）。convex combination 的结论要求 weather 无 skewness，且偏离平方要有 variation，所以这里把两县的偏离都换成对称的三点分布：$-2$、$0$、$+2\,^\circ$C，概率 $\tfrac14$、$\tfrac12$、$\tfrac14$。每县 $\sigma^2=2$，$\E[d^4]=8$，$v=8-2^2=4$。于是 $S=4$、$V=8$、$\lambda_{\text{甲}}=\lambda_{\text{乙}}=\tfrac12$、$\bar\mu=20\,^\circ$C；两县 climate 相差 $8\,^\circ$C，$\sum_{i<j}(\mu_i-\mu_j)^2\sigma_i^2\sigma_j^2=64\times2\times2=256$。代入式~\ref{eq:cca-theta}：
\begin{equation*}
\bar\theta=\frac{4\times8}{4\times8+4\times256}=\frac{32}{1056}=\frac{1}{33}\approx0.03 .
\end{equation*}
naive 二次面板的系数有约 97\% 的权重落在 long-run 系数上。取 $\beta_3=-0.01$：曲率的偏差是 $\beta_3\bar\theta\approx-0.0003$；marginal effect 的偏差在甲县的 climate（$\mu=16$）处是 $2\times(-0.01)\times\tfrac{1}{33}\times(16-20)\approx+0.0024$，在乙县（$\mu=24$）处是 $-0.0024$，在 $\bar\mu=20$ 处为零。再把两县的 climate 拉近到 $19\,^\circ$C 与 $21\,^\circ$C、weather 分布不变：成对差异项变成 $4\times2\times2=16$，$\bar\theta=32/(32+64)=\tfrac13$，short-run 的权重升到三分之一。同样的年际波动，横截面 climatic differences 越小，naive 系数里的 short-run 成分越多。这正是下文 Table~2 里温度与降水权重悬殊的原因。

\subsection{两个极限与适用边界（式 (13)、附录 F）}

式~\ref{eq:cca-theta} 的 $\bar\theta$ 何时接近零（即估计基本反映 long-run response）？论文的条件 (13) 是：分子那一项远小于成对 climatic differences 那一项，即各 location 内偏离平方的时间序列 variation 相对横截面 climatic differences 要小（第~1219 页）。一个极限情形值得单记：``if the time-series variation in the weather data is purely binary (say either hot or cold weather), then this simple source of variation allows parameter identification without causing bias''\citep[第~1219 页]{merel2021climate}。准确的条件是 weather 只取两个\textbf{等概率}的值（第~1216、1219 页）：此时偏离平方在 location 内是常数，$v_i=0$，penalty term 被 fixed effects 完全吸收，$\bar\theta=0$。等概率不可省略：上一小节线性例子里的甲县也只有两个取值，但概率是 0.8 与 0.2，分布右偏，偏差不为零。

第二点关于 climate 有趋势的情形（附录 F，正文第~1222 页转述）。若各 location 的 climate 带趋势，漏掉的 penalty term 不随 $T$ 增长，而回归元的 variation 随 $T$ 增长，naive 估计量相对 $\boldsymbol{\beta}^{\mathrm{LR}}$ 的偏差随 $T\to\infty$ 消失，论文称之为 \index{superconsistency}superconsistency（超一致收敛）；这是有趋势时才有的结论，stationary 情形下的式~\ref{eq:cca-combo} 本身已是 $T\to\infty$ 的极限，偏差并不消失。作者还指出一个反直觉之处：即使回归里给每个 location 加了时间趋势，superconsistency 依然成立，因为``detrending the squared variable is conceptually different than squaring the detrended weather variable''\citep[第~1222 页]{merel2021climate}。有趋势的模型在实际数据上的代价是：``the weight parameters obtained from the non-stationary models often lie outside the unit interval''\citep[第~1222 页]{merel2021climate}；而在美国、Corn Belt 与法国的数据上，stationarity assumption 下的式 (11)(12) 对估计响应曲线位置的预测``remarkably well''\citep[第~1222 页]{merel2021climate}——这是结果节采用 stationary 框架的经验依据。

\section{一致估计：climate penalty 增广回归}
\label{sec:cca-augmented}

偏差分解是诊断，论文的第二件贡献是给出对策：一个能同时给出 short-run response、long-run response 与 adaptation 检验的设定——``a parsimonious panel estimation framework that simultaneously delivers short-run responses to weather conditional on climate, a global long-run response to climate, and a simple yet formal test of climatic adaptation''\citep[第~1210 页]{merel2021climate}。

\subsection{增广设定与识别（式 (16)）}

\begin{ident}[climate penalty 增广回归：可直接搬用的设定]
\begin{align*}
\text{估计方程：}\quad & y_{it} = a_i + b_1 x_{it} + b_2 x_{it}^2 + b_3\,(x_{it}-\bar{x}_i)^2 + e_{it},\qquad \bar{x}_i\equiv\tfrac{1}{T}\textstyle\sum_t x_{it} \\
\text{long-run response：}\quad & b_1 x + b_2 x^2\ \text{（各 location 共有，差一个截距）} \\
\text{short-run response：}\quad & \text{上式加 penalty } b_3(x-\bar{x}_i)^2 \\
\text{identification：}\quad & \text{stationarity assumption 下 within estimator 一致（附录 C）}
\end{align*}
配套口径：$\bar{x}_i$ 取样本 weather 均值；可含时间效应等额外控制而不破坏一致性；可纳入多个 weather 变量，各自带自己的 penalty term（第~1220 页）。
\end{ident}

设定与式~\ref{eq:cca-dgp} 只差一处：未知的 $\mu_i$ 换成样本均值 $\bar{x}_i$。这并不破坏一致性：``We show in Appendix C that the within estimator of model (16) is consistent for $\boldsymbol{\beta}$ under the stationarity assumption. Thus, equation (16) yields consistent estimates of the long-run climate response as well as the adaptation parameter $\beta_3$, allowing one to gauge the extent to which long-run inputs contribute to the overall response to climate''\citep[第~1220 页；兑现第~1214 页脚注~17「正确设定要求知道 $\mu_i$，但一致估计不要求」的预告]{merel2021climate}。识别直觉由 Figure~4（第~1220--1221 页）给出：两个 location 在同一时期经历「对各自 climate 的相同绝对偏离」$A$，则 penalty $B$ 相同；``conditional on the penalty B, the estimated effect $\hat{b}_1 x + \hat{b}_2 x^2$ delineates the underlying long-run relationship, which is thus identified by comparing outcomes across observations with identical deviations but differing weather realizations''\citep[第~1221 页]{merel2021climate}——penalty 被条件化之后，剩下的二次函数就是各 location 共有的 long-run response，short-run response 再叠加各自的 penalty。

\subsection{counterfactual 计算（式 (17)）与三条注意事项}

设 location 的 climate 从 $\mathbf{C}_0=(\mu_0,\sigma_0^2)$ 变到 $\mathbf{C}_1=(\mu_1,\sigma_1^2)$，均值与方差都可以变（论文各量带 location 下标 $i$，此处从略）。在 climate $k$ 下期望产出是 $\alpha_i+\beta_1\mu_k+\beta_2(\mu_k^2+\sigma_k^2)+\beta_3\sigma_k^2$，于是长期 counterfactual 为（第~1221 页）
\begin{equation*}
\Delta^{\mathrm{LR}}\widehat{\E}[y] = \widehat{b}_1(\mu_1-\mu_0) + \widehat{b}_2(\mu_1^2-\mu_0^2) + (\widehat{b}_2+\widehat{b}_3)(\sigma_1^2-\sigma_0^2),
\end{equation*}
短期 counterfactual （行动仍按旧 climate 安排）为
\begin{align}
\Delta^{\mathrm{SR}}\widehat{\E}[y] \;=\;& \widehat{b}_1(\mu_1-\mu_0) + \widehat{b}_2\bigl(\mu_1^2-\mu_0^2+\sigma_1^2-\sigma_0^2\bigr) \notag\\
&+ \widehat{b}_3\bigl(\sigma_1^2-\sigma_0^2+(\mu_1-\mu_0)^2\bigr)
\;=\; \Delta^{\mathrm{LR}}\widehat{\E}[y] + \widehat{b}_3(\mu_1-\mu_0)^2 .
\label{eq:cca-cf}
\end{align}
（式 (17)，第~1221 页；$\sigma_k^2$ 为 climate $k$ 下 weather 的方差。）两者只差一项 $\widehat{b}_3(\mu_1-\mu_0)^2$：climate 均值移动之后、行动尚未调整时多付的 penalty。论文对这组公式给了三条 remarks。其一，长期期望产出仍含 $\beta_3\sigma^2$ 项，原因是``agents adapt to climate but still face year-to-year variation in weather, hence the presence of the term in $\beta_3$ in the long-run expected outcome. (Long-run adaptation means that agents move from one short-run response curve to a different one when climate changes, not that yearly outcomes occur on the long-run response function.)''\citep[第~1221 页]{merel2021climate}。其二，若 counterfactual 不涉及变率的变化（$\sigma_1^2=\sigma_0^2$），$\widehat{b}_3$ 从长期 counterfactual 中消去，只用 long-run response function 即可（第~1221--1222 页）。其三，若变率有变化而仍按 naive 路线计算（相当于令 $\widehat{b}_3=0$），预测有偏；不少地区的 climate 模型预测温度或降水的年际变率上升，而 $\beta_3\le 0$，故这样的预测偏乐观（第~1222 页）。

\section{数据与测度}

论文的经验部分用三组已发表研究的面板（作者声明``we revisit panel estimates from existing studies''\citep[第~1224 页]{merel2021climate}）。美国的 weather 数据由 Wolfram Schlenker 提供（脚注~26），国家面板取自 Burke 等的斯坦福数据页（脚注~27，均在第~1224 页）。

\textbf{美国县面板。}66 年的面板，变量为 growing season（4--10 月，玉米与大豆）的平均温度、growing degree days（定义为 10--29\,$^\circ$C）与总降水，覆盖 100 度经线以东、不含 Florida 的县，$I=2{,}391$（第~1224 页）。子样本为 Corn Belt 三州 Iowa、Illinois、Indiana，$I=293$（第~1224 页）。产出是对数单产，数据来自 USDA NASS（第~1225 页脚注~28）；控制 county 与 state$\times$year fixed effects，标准误在州年层面聚类（第~1225--1226 页）；另有一组替代设定，以 growing degree days 的二次式刻画温度暴露，并用 heating degree days 的平方根控制极端高温，沿 Schlenker、Hanemann 与 Fisher（2007）与 Schlenker（2017）（第~1225 页）。

\textbf{法国省面板。}66 年的法国 department 级（省级）面板，$I=88$，变量为小麦与大麦 growing season（3--7 月）的平均温度与累积降水；该数据此前在 Gammans、Mérel 与 Ortiz-Bobea（2017）中用过（第~1224 页）。

\textbf{全球国家面板。}Burke、Hsiang 与 Miguel（2015）与 Burke、Davis 与 Diffenbaugh（2018）的国家年度均温与降水面板，$I=167$，结果在附录 G.4（第~1224--1225 页）。

\begin{pitfall}[权重参数只讲 variation 结构，不讲 adaptation 是否发生]
Table~2 的 $\bar\theta$（与 $\tilde\theta$）只依赖历史 weather 数据本身：``Importantly, these values depend only on the historical weather data and are independent of the outcome of interest, reflecting the ability of a given weather dataset to capture adaptive behavior if present, but they do not indicate whether adaptation has occurred in a given empirical setting''\citep[第~1225 页]{merel2021climate}。误读形态：把「温度的 $\bar\theta$ 只有 2\%」当成「美国玉米对温度的 adaptation 很强」。前者说的是这套数据的 variation 绝大部分来自横截面（所以 naive 系数接近 long-run response），后者要靠增广回归里 $\widehat{\beta}_3$ 的显著性与量级来回答（Table~3--4）。
\end{pitfall}

\section{结果与解释}

\subsection{第一步：variation 构成（Table 1）}

论文在「Trends in Historical Climate」一节先核对 stationary 框架是否够用。Table~1（第~1223 页；对比样本 $T=46$，1970--2015 年）报告两组份额。第一组是各 location 的 weather 时间 variation 中来自逐年波动（而非 climate 趋势）的份额：温度为美国 0.914、Corn Belt 0.959、法国 0.463，降水为 0.977、0.963、0.982。第二组是横截面 climatic variation 中独立于 climate 趋势的份额：温度为 0.996、0.987、0.797，降水为 0.955、0.611、0.970。读法：除法国的 growing season 温度（趋势占时间 variation 的份额不可忽略）之外，climate 趋势只解释 location 内 weather variation 的很小部分；美国样本中超过 99\% 的 climatic temperature variation 与超过 95\% 的 climatic precipitation variation 独立于趋势，法国为 80\% 与 97\%，Corn Belt 降水只有 61\%，原因是该样本降水的横截面 variation 本来就有限（第~1223--1224 页）。结论：``in our settings, climatic trends tend to be small relative to both year-to-year weather fluctuations and persistent cross-sectional climatic differences''\citep[第~1224 页]{merel2021climate}——stationary 框架在本章数据上大体够用。

\subsection{第二步：权重参数（Table 2）}

Table~2（第~1225 页）给出各面板各变量的 $\bar\theta$（无 skewness 的式 (12)）、$\tilde\theta$（容许 skewness 的式 (15)）以及对应的组合 climate $\bar\mu$ 与 $\tilde\mu$；括号内是按年份有放回重抽 1{,}000 次得到的标准差；$\tilde\mu$ 落在样本 climate 的范围之内，通常接近 $\bar\mu$（第~1225 页）。温度的 $\bar\theta$ 极小：美国县面板 0.022、Corn Belt 0.123、法国 0.141——即 long-run 的权重分别约 98\%、88\%、86\%：``For temperature data, we find that estimates of the effect of growing-season temperature should mostly reflect the long-run response''\citep[第~1225 页]{merel2021climate}。降水的权重则大得多且跨数据集差异悬殊：``For the US, Corn Belt, and French panels, the weight on the long-run response is 41\%, 7\%, and 53\%, respectively''\citep[第~1225 页]{merel2021climate}。作者提醒两点：$\bar\theta$ 只讲 variation 结构（见上节易错点）；容许 skewness 的 $\tilde\theta$ 与 $\bar\theta$ 相差不大（温度的 $\tilde\theta$ 为 0.023、0.124、0.146），``allowing for skewness in the weather data does not meaningfully change the calculated values''\citep[第~1225 页]{merel2021climate}，所以论文的讨论以 $\bar\theta$ 为准。

\subsection{第三步：增广回归的 adaptation 检验（Table 3--4）}

把 penalty term 放进回归，$\widehat{\beta}_3$ 的显著性就是 climatic adaptation 的检验。读表先看单位：温度 penalty 系数乘以 100 报告，降水 penalty 系数乘以 $10^6$ 报告。美国全县样本（Table~3，第~1226 页）：玉米的温度 penalty 为 $-3.87$（标准误 1.31）、降水 penalty 为 $-0.37$（0.12），均在 1\% 水平显著；``We find significant adaptation to both temperature and precipitation, as evidenced by the fact that short-run responses lie below long-run responses, and the difference in concavity between the curves, which is entirely due to the adaptation coefficient $\beta_3$, is statistically significant''\citep[第~1226--1227 页]{merel2021climate}。大豆的温度 penalty 为 $-2.20$（0.74），在 1\% 水平显著，降水 penalty 为 $-0.22$（0.13），只在 10\% 水平显著；量级都小于玉米，``suggesting more limited scope for climatic adaptation in soybeans''\citep[第~1230 页]{merel2021climate}。Corn Belt 子样本（Table~4，第~1226 页）：玉米的温度 penalty 为 $-5.72$（2.06），在 1\% 水平显著，降水 penalty 为 $-0.54$（0.57），不显著（第~1230 页）；大豆的温度 penalty 为 $-2.29$（1.26），只在 10\% 水平显著，降水 penalty 为 $-0.69$（0.43），不显著（第~1231 页）。图形版本见论文 Figures~5--8（第~1227--1230 页）。以全样本玉米为例（第~1227 页）：对温度的 long-run adaptation 几乎完全反映在 naive 估计的响应里，因为横截面 variation 占主导；对降水，naive 响应明显比 long-run 响应更凹，weather 波动在面板总 variation 中的作用更大。这与第二步的权重是一致的。

\subsection{第四步：counterfactual 量化}

按式~\ref{eq:cca-cf} 的口径，论文报告两套数。第一套是在单个（虚构的）location 上评价的 counterfactual，针对美国玉米：2\,$^\circ$C 增温的 short-run 影响为 $-41\%$、long-run 为 $-26\%$，``Thus, long-run adaptation buffers yield losses by more than fifteen percentage points''\citep[第~1228 页]{merel2021climate}；降水减 20\% 为 short-run $-4.6\%$ 对 long-run $-3.9\%$，增 20\% 为 short-run $+1.4\%$ 对 long-run $+2.1\%$（第~1229 页）。第二套是按面积权重加总的总量 counterfactual：2\,$^\circ$C 增温使美国玉米单产 short-run $-27\%$、long-run $-16\%$，大豆 $-21\%$ 对 $-15\%$，法国小麦 $-10\%$ 对 $-7\%$（引言，第~1208 页；正文分别在第~1228、1230、1231 页）。两套数都出自增广回归，差别只在评价对象：论文说明第一套数只对应单个虚构的 location，不代表全国产出的总效应（第~1228 页）。\textbf{引用时先分清口径}：单点数与加总数不得混用。

\subsection{第五步：遗留偏差与全球面板}

论文在经验部分末尾专设「Vulnerability of Panel Estimates to Omitted Variables」一小节（第~1233--1235 页），指出增广回归仍可能有偏的一个来源：不随时间变化、又与 weather 交互的未观测变量（式 (18)(19)）。作者提出两条处理思路（第~1235 页）：一是让式 (16) 中的 $b_1$ 随 location 变化，以非参数的方式控制未观测的交互，同时允许 climate 随时间变化以保住识别，这要求各地 climate 有足够的趋势；二是 Druckenmiller 与 Hsiang（2018）的 spatial first differences，在 climate 趋势很小或不存在的情形下可作替代。全球国家面板的结果在附录 G.4（人均 GDP 增长，第~1225 页），本仓未收，待补。stationary 框架是否够用，正文的依据是附录 F.4 在美国、Corn Belt 与法国数据上的比较：有趋势的模型给出的权重参数常落在单位区间之外，而 stationarity assumption 下的式 (11)(12) 对估计响应曲线位置的预测``remarkably well''\citep[第~1222 页]{merel2021climate}。

\section{局限与批评}

论文自己承认的：其一，二次设定之下的结论有层次之分——一般情形下偏差的符号不能确定，也不再完全由 skewness 驱动（第~1216--1217 页）；convex combination 与 $\bar\theta$ 的干净表达式依赖无 skewness 与 stationary climate 两个简化（「渐近偏差分解」节已随文标注）。其二，stationary climate assumption：放松后权重参数可能落到单位区间外（第~1222 页）。其三，counterfactual 对变率变化的处理：忽略变率变化的 naive 计算偏乐观（第~1222 页）。其四，增广回归对「与 weather 交互的未观测变量」并不免疫（第~1233--1235 页）；论文在引言里的总体态度是，即便横截面 climate variation 占主导，由非线性面板得出的 climate 影响也应谨慎解读（第~1210 页）。论文对增广设定的建议是积极的：它不比现行做法多要求任何信息，``we see no good reason not to adopt it in future work''\citep[第~1236 页]{merel2021climate}，并认为凡是结果变量对同一条件的短期与长期响应可以不同的面板研究都用得上（第~1236 页）。

后续文献的批评与发展：应用该框架的文章本讲义已读一篇，即 \chref{ch:da-et-al-2024}（中国冬小麦），它显示 penalty term 全部加入时会出现反常符号，以及 climate 有趋势的短面板上 stationarity 如何勉强维持；papers/01\_methods/ 下 2026 年 AJAE 一篇（加州多年生作物）尚未读，待补写。

本讲义的评注：这篇论文与 \chref{ch:hsiang-2016} 应对照读——Hsiang 的 Envelope 论证说「在已优化的 climate 处，weather 的 marginal effect 等于 climate 的 marginal effect」，论文的回答是「漏掉 penalty term 的面板估计到的是加权平均，只在组合 climate $\bar{\mu}$ 处 marginal effect 才无偏」（``our result is both different and more specific''\citep[第~1217 页]{merel2021climate}）。方法线上，D-G 的自认局限（\chref{ch:deschenes-greenstone-2007}，原文第~359 页）在此第一次有了可计算的答案，BHM 的系数解读问题（\chref{ch:burke-hsiang-miguel-2015}）在此有了量化工具。

\section{在本领域脉络中的位置}

论文回答的旧问题：panel approach 与 nonlinear panel approach 的系数各测到什么（proposition 的「多少」）。它用的路线：McIntosh 与 Schlenker（2006）的 DGP（\citealp{mcintosh2006identifying}）、Burke 与 Emerick（2016）式的 short-run 对 long-run 比较（作为隐含检验重现于同一回归，第~1210 页）、Hsiang（2016）的边际等价论证（被明确讨论并限定，第~1217 页）。它开的新问题：宏观加总损害函数的构造——弄清损害估计在多大程度上反映 long-run 值，``is imperative for the construction of macro-level aggregate damage functions, which in turn are necessary for calculating the social cost of carbon''\citep[第~1235 页]{merel2021climate}；增广框架在后续农业研究中的推广，已有一例见 \chref{ch:da-et-al-2024}。本章在 \chref{ch:info-climate-change-adaptation} 的方法线（panel 线第三站）、发现线与速查表中均有登记。

% 体例备注（2026-11-04）：「机制与讨论」按体例属视需要小节，本篇的机制内容
% （penalty 行为推导、tangency 几何、adaptation 检验）已并入第二节与
% 「一致估计」「结果与解释」各节，不另设。

% 章末参考文献（只列本章引用的条目），勿删
\chapterrefs

% ---- 应用篇：把上面的方法用到具体作物／地区的各章，排在方法篇之后，
%      按 \include 顺序编号 Application 1、2……（每个 part 各自从 1 起）。
%      编号不进章题（章题＝论文标题），写在四处：\chapter 的可选参数（目录条目）、
%      \chaptermark（页眉）、该章「它在 partinfo 里的坐标」一句、partinfo 速查表。
%      调整顺序或插入新的应用篇时，这四处的编号要一起改。 ----
\chapter{Application1: Climate change and winter wheat yields: the role of freezing days and long-run adaptations}

\timestamp{2026-11-11}{2026-11-11}

\section{这篇文章在说什么}
本论文在估计 climate change impacts on Chinese winter wheat yields的过程中使用了season-specific climate variables和climate penalty mode，得出了两个结论：1、在实证估计中纳入季节性变量是必要的，本研究中Spring freezing days的effects与传统估计完全不同；2、长期适应行为欸也会对climate change effect产生缩小甚至逆转的效果。

\section{数据来源}
\subsection{数据规模：region \times period}
\subsection{yield data}
\subsection{climate data}

\section{temperature thresholds estimate}
Step1: Seasonal division

Step2: Variables construction
figure 1

\begin{card}
    Question1: sum of four variables in given one season is constant，似乎会出现完全共线性问题？
    这就是为什么这里使用的是degree variables而不是bin variables，degree variables可以避免完全共线性问题。
\end{card}

Step3: Range loop
table 1

\begin{card}
    "pick the one which has the best fit"

    What's the fit function? 论文正文中并没有，从Appendix或者reference中寻找补充到此处。
\end{card}

\section{The estimate model}
\subsection{The conventional fixed effects panel data model}
equation (1) in paper，有seasonal variables，只是没有添加climate penalty terms

\subsection{The climate penalty model}
A simple climate penalty model: equation (2) in paper

how to interpret coefficients: \beta_1 代表 long-response， (\beta_1, \beta_2)代表 short-response。

\begin{card}
    Trending Historical change，将Appendix中的内容补充到这里。
\end{card}

\section{Empirical results}
Table 2
\subsection{Column 1: Conventional panel approach}
解释系数，并强调这些系数是long-response.

\subsection{COlumn 2/3: Panel approaches with climate penalty terms}
解释系数，
contrast with Column1
Interesting findings and interprations

\subsection{Robustness check}

\subsection{Projecting}
包含 5.4 和5.5 部分。                     % Application 1

% 书末总参考文献：先关闭按章分单元（bibunits 手册 §3.3），
% globalcitecopy 选项已把各章引用汇入此表
\bibliographyunit
\bibliographystyle{plainnat}
\bibliography{bibs/refs}

\end{document}
